\documentclass[11pt,letterpaper]{article}


% pmml  arff  openannotation

%\usepackage[condensed,math]{anttor}
%\usepackage[T1]{fontenc}

%\usepackage[T1]{fontenc}
%\usepackage{tgtermes}

\usepackage[hang,flushmargin]{footmisc}

\usepackage{titlesec}

%\usepackage{sectsty}
%\sectionfont{\fontsize{13}{4}\selectfont}

\titleformat{\section}
  {\normalfont\fontsize{13}{15}\bfseries}{\thesection}{1em}{}

\let\OldPart\part

\renewcommand{\part}[1]{\OldPart{#1}%
%{\textcolor{darkRed}\hrule}
\vspace{-.5em}}

\titlespacing*{\section}
{0pt}{4ex plus 1ex minus .1ex}{-0.2ex plus .2ex}

\titlespacing*{\subsection}
{0pt}{3.5ex plus 1ex minus .1ex}{-.8ex plus .2ex}

%\usepackage{mathptmx}

\usepackage{eso-pic}

%\setlength\parindent{0pt}

\AddToShipoutPictureBG{%

\ifnum\value{page}>1{
\AtTextUpperLeft{
\makebox[20.5cm][r]{
\raisebox{-1.95cm}{%
{\transparent{0.3}{\includegraphics[width=0.29\textwidth]{e-logo.png}}	}} } }
}\fi
}

\AddToShipoutPicture{%
{
 {\color{blGreen!70!red}\transparent{0.9}{\put(0,0){\rule{3pt}{\paperheight}}}}%
 {\color{darkRed!70!purple}\transparent{1}\put(3,0){{\rule{4pt}{\paperheight}}}}
% {\color{logoPeach!80!cyan}\transparent{0.5}{\put(0,700){\rule{1cm}{.6cm}}}}%
% {\color{darkRed!60!cyan}\transparent{0.7}\put(0,706){{\rule{1cm}{.6cm}}}}
% \put(18,726){\thepage}
% \transparent{0.8}
}
}

\AddToShipoutPicture{%
\ifnum\value{page}=1
\put(257.5,942){%
	\transparent{0.7}{
		\includegraphics[width=0.2\textwidth]{logo.png}}}
\put(59,953){\textbf{{\fontfamily{phv}\fontsize{14}{14}\selectfont{}WHITE PAPER}}}
\fi
}	



\AddToShipoutPicture{%
\ifnum\value{page}>1
{\color{blGreen!70!red}\transparent{0.9}{\put(300,8){\rule{0.5\paperwidth}{.3cm}}}}%
{\color{inOne}\transparent{0.8}{\put(300,10){\rule{0.5\paperwidth}{.3cm}}}}%
{\color{inTwo}\transparent{0.3}\put(300,13){{\rule{0.5\paperwidth}{.3cm}}}}

%\ifnum\value{page}<39
\put(301,16){%
\transparent{0.7}{
\includegraphics[width=0.2\textwidth]{logo.png}}
}
%\fi

%\pgfmathparse{equal(\value{page},15)||equal(\value{page},20)?int(1):int(0)}

\ifnum\value{page}=8   %\pgfmathresult>0\relax
\put(182,142){
 {\setlength{\fboxsep}{.65em}\fontsize{9}{10}\selectfont

   {\color{white}{\parbox{11cm}{\vspace{7pt}\framebox{\begin{minipage}{.46\textwidth}
	  {\color{black}{\textit{For more information please contact:}}\\\textbf{\color{blGreen!40!blbl}{Amy Neustein, Ph.D., Founder and CEO}}\\  
       \textbf{{\color{blGreen!20!black}Linguistic Technology Systems \\
      amy.neustein@verizon.net \textbullet{} \textbf{(917) 817-2184} }}}
	 \end{minipage}}}}}}
}
\fi

{\color{blGreen!70!red}\transparent{0.9}{\put(5.6,5){\rule{0.5\paperwidth}{.4cm}}}}%
{\color{inOne}\transparent{1}{\put(5.6,10){\rule{0.5\paperwidth}{.4cm}}}}%
{\color{inTwo}\transparent{0.3}\put(5.6,15){{\rule{0.5\paperwidth}{.4cm}}}}

\fi
}

%\pagestyle{empty} % no page number
%\parskip 7.2pt    % space between paragraphs
%\parindent 12pt   % indent for new paragraph
%\textwidth 4.5in  % width of text
%\columnsep 0.8in  % separation between columns

%\setlength{\footskip}{7pt}

\usepackage[paperheight=14in,paperwidth=8.5in]{geometry}
\geometry{left=.83in,top=.45in,right=.81in,bottom=1.2in} %margins

\usepackage{etoolbox}% http://ctan.org/pkg/etoolbox
\makeatletter
% \patchcmd{<cmd>}{<search>}{<replace>}{<success>}{<failure>}
\patchcmd{\@part}{\par}{\quad}{}{}
\patchcmd{\@part}{\huge}{\Large}{}{}
\makeatother

\renewcommand{\partname}{\hspace{-1em}Part}

\renewcommand*\thepart{\Roman{part}:}

\renewcommand{\thepage}{\raisebox{2pt}{\arabic{page}}}

\renewcommand{\footnoterule}{%
	\kern 4pt
	\hrule width .92\textwidth height .5pt
	\kern 6pt
}


\usepackage[hyphens]{url}
\newcommand{\biburl}[1]{ {\fontfamily{gar}\selectfont{\textcolor[rgb]{.2,.6,0}%
{\scriptsize {\url{#1}}}}}}

%\linespread{1.3}

\newcommand{\sectsp}{\vspace{12pt}}

\usepackage{graphicx}
\usepackage{color,framed}

\usepackage{textcomp}

\usepackage{float}

\usepackage{mdframed}


\usepackage{setspace}
\newcommand{\rpdfNotice}[1]{\begin{onehalfspacing}{

\Large #1

}\end{onehalfspacing}}

\usepackage{xcolor}

\usepackage[hyphenbreaks]{breakurl}
\usepackage[hyphens]{url}

\usepackage{hyperref}
\newcommand{\rpdfLink}[1]{\href{#1}{\small{#1}}}
\newcommand{\dblHref}[1]{\href{#1}{\small{\burl{#1}}}}
\newcommand{\browseHref}[2]{\href{#1}{\Large #2}}

\colorlet{blCyan}{cyan!50!blue}

\definecolor{darkRed}{rgb}{.2,.0,.1}


\definecolor{blGreen}{rgb}{.2,.7,.3}

\definecolor{darkBlGreen}{rgb}{.1,.3,.2}

\definecolor{oldBlColor}{rgb}{.2,.7,.3}

\definecolor{blColor}{rgb}{.1,.3,.2}

\definecolor{elColor}{rgb}{.2,.1,0}
\definecolor{flColor}{rgb}{0.7,0.3,0.3}

\definecolor{logoOrange}{RGB}{108, 18, 30}
\definecolor{logoGreen}{RGB}{85, 153, 89}
\definecolor{logoPurple}{RGB}{200, 208, 30}

\definecolor{logoBlue}{RGB}{4, 2, 25}
\definecolor{logoPeach}{RGB}{255, 159, 102}
\definecolor{logoCyan}{RGB}{66, 206, 244}
\definecolor{logoRed}{rgb}{.3,0,0}

\newcommand{\colorq}[1]{{\color{logoOrange!70!black}{\q{\small\textbf{#1}}}}}

\definecolor{inOne}{rgb}{0.122, 0.435, 0.698}% Rule colour
\definecolor{inTwo}{rgb}{0.122, 0.698, 0.435}% Rule colour

\definecolor{outOne}{rgb}{0.435, 0.698, 0.122}% Rule colour
\definecolor{outTwo}{rgb}{0.698, 0.435, 0.122}% Rule colour

\colorlet{linkcolor}{flColor!60!red}


\hypersetup{
	colorlinks=true,
	citecolor=blCyan!40!green,
	filecolor=magenta!30!logoBlue,
	urlcolor=blue,
    linkcolor=linkcolor!70!black,
%    allcolors=blCyan!40!green
}


\usepackage[many]{tcolorbox}% http://ctan.org/pkg/tcolorbox

\usepackage{transparent}

\newlength{\bsep}
\setlength{\bsep}{-1pt}
\let\xbibitem\bibitem
\renewcommand{\bibitem}[2]{\vspace{\bsep}\xbibitem{#1}{#2}}

\newenvironment{cframed}{\begin{mdframed}[linecolor=logoPeach,linewidth=0.4mm]}{\end{mdframed}}

\newenvironment{ccframed}{\begin{mdframed}[backgroundcolor=logoGreen!5,linecolor=logoCyan!50!black,linewidth=0.4mm]}{\end{mdframed}}


%\usepackage[T1]{fontenc}

%\usepackage{aurical}
% \Fontauri

\usepackage{gfsdidot}
\usepackage[T1]{fontenc}

%\makeatletter
%\f@family,  cmr, T1, n, m,
%\f@encoding,
%\f@shape,
%\f@series,
%\makeatother



%\usepackage{LibreBodoni}

%\usepackage{fontspec}
%\setmainfont{QTBengal}

\usepackage{relsize}

\newcommand{\bref}[1]{\hspace*{1pt}\textbf{\ref{#1}}}

\newcommand{\pseudoIndent}{

\vspace{10pt}\hspace*{12pt}}

\newcommand{\YPDFI}{{\fontfamily{fvs}\selectfont YPDF-Interactive}}

%
\newcommand{\deconum}[1]{{\protect\raisebox{-1pt}{{\LARGE #1}}}}

\newcommand{\visavis}{vis-\`a-vis}

\newcommand{\VersatileUX}{{\color{red!85!black}{\Fontauri Versatile}}%
{{\fontfamily{qhv}\selectfont\smaller UX}}}

\newcommand{\NDPCloud}{{\color{red!15!black}%
{\fontfamily{qhv}\selectfont {\smaller NDP C{\smaller LOUD}}}}}

\newcommand{\MThreeK}{{\color{blGreen!45!black}%
{\fontfamily{qhv}\fontsize{10}{8}\selectfont {M3K}}}}


\newcommand{\lfNDPCloud}{{\color{red!15!black}%
{\fontfamily{qhv}\selectfont N{\smaller DP C{\smaller LOUD}}}}}

\newcommand{\textds}[1]{{\fontfamily{lmdh}\selectfont{%
\raisebox{-1pt}{#1}}}}

\newcommand{\dsC}{{\textds{ds}{\fontfamily{qhv}\selectfont \raisebox{-1pt}
{\color{red!15!black}{C}}}}}

\definecolor{tcolor}{RGB}{24,52,61}

\newcommand{\CCpp}{\resizebox{!}{7pt}{\AcronymText{C}}/\Cpp{}}
\newcommand{\NoSQL}{\resizebox{!}{7pt}{\AcronymText{NoSQL}}}
\newcommand{\SQL}{\resizebox{!}{7pt}{\AcronymText{SQL}}}

\newcommand{\SPARQL}{\resizebox{!}{7pt}{\AcronymText{SPARQL}}}

\newcommand{\NCBI}{\resizebox{!}{7pt}{\AcronymText{NCBI}}}

\newcommand{\HTXN}{\resizebox{!}{7pt}{\ATexttclr{HTXN}}}

\newcommand{\TDM}{\resizebox{!}{7pt}{\AcronymText{TDM}}}

\newcommand{\lHTXN}{\resizebox{!}{7.5pt}{\ATexttclr{H}}%
\resizebox{!}{6.5pt}{\ATexttclr{TXN}}}

\newcommand{\lsHTXN}{\resizebox{!}{9.5pt}{\ATexttclr{HTXN}}}

\newcommand{\LAF}{\resizebox{!}{7pt}{\AcronymText{LAF}}}

\newcommand{\UDpipe}{\resizebox{!}{7pt}{\AcronymText{UDpipe}}}

\newcommand{\C}{\resizebox{!}{7pt}{\AcronymText{C}}}

\newcommand{\FCS}{\resizebox{!}{7pt}{\AcronymText{FCS}}}

\newcommand{\GAVI}{\resizebox{!}{7pt}{\AcronymText{GAVI}}}
\newcommand{\ArcGIS}{\resizebox{!}{7pt}{\AcronymText{ArcGIS}}}
\newcommand{\QGIS}{\resizebox{!}{7pt}{\AcronymText{QGIS}}}

\newcommand{\GIS}{\resizebox{!}{7pt}{\AcronymText{GIS}}}
\newcommand{\AngelScript}{\resizebox{!}{7pt}{\AcronymText{AngelScript}}}



\usepackage{mdframed}

\newcommand{\cframedboxpanda}[1]{\begin{mdframed}[linecolor=yellow!70!blue,linewidth=0.4mm]#1\end{mdframed}}


\newcommand{\PVD}{\resizebox{!}{7pt}{\AcronymText{PVD}}}

\newcommand{\TAGML}{\resizebox{!}{7pt}{\AcronymText{TAGML}}}
\newcommand{\GTagML}{\resizebox{!}{7pt}{\AcronymText{GTagML}}}

\newcommand{\lGTagML}{\resizebox{!}{7.5pt}{\AcronymText{GTagML}}}




\newcommand{\SDK}{\resizebox{!}{7pt}{\AcronymText{SDK}}}
\newcommand{\NLP}{\resizebox{!}{7pt}{\AcronymText{NLP}}}

\newcommand{\AXF}{\resizebox{!}{7pt}{\ATexttclr{AXF}}}

\newcommand{\HyperGraphDB}{\resizebox{!}{7pt}{\AcronymText{HyperGraphDB}}}

\newcommand{\AllegroGraph}{\resizebox{!}{7pt}{\AcronymText{AllegroGraph}}}

\newcommand{\Grakenai}{\resizebox{!}{7pt}{\AcronymText{Graken.ai}}}


\newcommand{\lAXF}{\resizebox{!}{7.5pt}{\ATexttclr{A}}%
\resizebox{!}{6.5pt}{\ATexttclr{XF}}}


\newcommand{\lsAXF}{\resizebox{!}{8.5pt}{\ATexttclr{AXF}}}

\newcommand{\AXFD}{\resizebox{!}{7pt}{\ATexttclr{AXFD}}}

\newcommand{\CBICA}{\resizebox{!}{7pt}{\AcronymText{CBICA}}}

\newcommand{\IORT}{\resizebox{!}{7pt}{\AcronymText{IORT}}}

\newcommand{\openCyto}{\resizebox{!}{7pt}{\AcronymText{openCyto}}}
\newcommand{\FACSanadu}{\resizebox{!}{7pt}{\AcronymText{FACSanadu}}}



\newcommand{\SeDI}{\resizebox{!}{7pt}{\AcronymText{SeDI}}}
\newcommand{\RSNA}{\resizebox{!}{7pt}{\AcronymText{RSNA}}}

\newcommand{\CER}{\resizebox{!}{7pt}{\AcronymText{CER}}}
\newcommand{\PACS}{\resizebox{!}{7pt}{\AcronymText{PACS}}}

\newcommand{\DCMTK}{\resizebox{!}{7pt}{\AcronymText{DCMTK}}}

\newcommand{\DICOM}{\resizebox{!}{7pt}{\AcronymText{DICOM}}}
\newcommand{\lsDICOM}{\resizebox{!}{9pt}{\AcronymText{DICOM}}}

\newcommand{\CytometryML}{\resizebox{!}{7pt}{\AcronymText{CytometryML}}}

\newcommand{\FCM}{\resizebox{!}{7pt}{\AcronymText{FCM}}}

\newcommand{\OMETIFF}{\resizebox{!}{7pt}{\AcronymText{OME-TIFF}}}
\newcommand{\OME}{\resizebox{!}{7pt}{\AcronymText{OME}}}


\newcommand{\CT}{\resizebox{!}{7pt}{\AcronymText{CT}}}

\newcommand{\LOINC}{\resizebox{!}{7pt}{\AcronymText{LOINC}}}

\newcommand{\RadLex}{\resizebox{!}{7pt}{\AcronymText{RadLex}}}


\newcommand{\OMOP}{\resizebox{!}{7pt}{\AcronymText{OMOP}}}
\newcommand{\PCORnet}{\resizebox{!}{7pt}{\AcronymText{PCORnet}}}
\newcommand{\FHIR}{\resizebox{!}{7pt}{\AcronymText{FHIR}}}

\newcommand{\CaPTk}{\resizebox{!}{7pt}{\AcronymText{CaPTk}}}

\newcommand{\VIOLIN}{\resizebox{!}{7pt}{\AcronymText{VIOLIN}}}



\newcommand{\lAXFD}{\resizebox{!}{7.5pt}{\ATexttclr{A}}%
\resizebox{!}{6.5pt}{\ATexttclr{XFD}}}


\newcommand{\IJST}{\resizebox{!}{7pt}{\AcronymText{IJST}}}

\newcommand{\BioC}{\resizebox{!}{7pt}{\AcronymText{BioC}}}

\newcommand{\CoNLL}{\resizebox{!}{7pt}{\AcronymText{CoNLL}}}
\newcommand{\CoNLLU}{\resizebox{!}{7pt}{\AcronymText{CoNLL-U}}}

\newcommand{\sapp}{\resizebox{!}{7pt}{\AcronymText{Sapien+}}}
\newcommand{\lsapp}{\resizebox{!}{8.5pt}{\AcronymText{Sapien+}}}
\newcommand{\lssapp}{\resizebox{!}{9.5pt}{\AcronymText{Sapien+}}}

\newcommand{\ePub}{\resizebox{!}{7pt}{\AcronymText{ePub}}}

%\lsLPF


\newcommand{\GIT}{\resizebox{!}{7pt}{\AcronymText{GIT}}}

%\definecolor{atColor}{RGB}{11, 71, 17}


\DeclareMathVersion{fordg}
\SetSymbolFont{letters}{fordg}{OML}{cmr}{b}{n}

\definecolor{atcColor}{RGB}{96, 17, 12}
%\textcolor{tcolor}{

\newcommand{\ATextCClr}[1]{\textcolor{atcColor}{\textbf{#1}}}

\newcommand{\ATexttclr}[1]{\textcolor{tcolor}{\textbf{#1}}}

\newcommand{\AIMConc}{\resizebox{!}{7.5pt}{\ATextCClr{AIM-Concepts}}}
\newcommand{\lAIMConc}{\resizebox{!}{8pt}{\ATextCClr{AIM-Concepts}}}

\newcommand{\HGXF}{{\resizebox{!}{7.5pt}{\ATexttclr{HGXF}}}}
\newcommand{\lHGXF}{{\resizebox{!}{8pt}{\ATexttclr{HGXF}}}}
\newcommand{\sHGXF}{{\resizebox{!}{6pt}{\ATexttclr{HGXF}}}}

\newcommand{\CRtwo}{{\resizebox{!}{7.5pt}{\ATextCClr{CR2}}}}
\newcommand{\lCRtwo}{{\resizebox{!}{8pt}{\ATextCClr{CR2}}}}
\newcommand{\sCRtwo}{{\resizebox{!}{6pt}{\ATextCClr{CR2}}}}


%\newcommand{\ConceptsDB}{{\resizebox{!}{7.5pt}{\ATextCClr{ConceptsDB}}}}

\newcommand{\THQL}{\resizebox{!}{7.5pt}{\ATexttclr{THQL}}}
\newcommand{\lTHQL}{\resizebox{!}{8pt}{\ATexttclr{THQL}}}

\newcommand{\HDICOM}{\resizebox{!}{7.5pt}{\ATexttclr{{\large h}-DICOM}}}

\newcommand{\hVaImm}{\resizebox{!}{7.5pt}{\ATexttclr{{\large h}-VaImm}}}


\newcommand{\PhaonVI}{\resizebox{!}{7.5pt}{\ATexttclr{Phaon-VI}}}



\definecolor{atColor}{RGB}{50, 22, 40}
\newcommand{\ATextClr}[1]{\textcolor{atColor}{\textbf{#1}}}

\newcommand{\DgDb}{{\mathversion{fordg}%
\makebox{\raisebox{-3pt}{\resizebox{!}{11pt}{\ATextClr{%
\rotatebox{17}{$\varsigma$}}}}\hspace{-4pt}%
\resizebox{!}{6.5pt}{\ATextClr{D\hspace{-2pt}B}}}}}


\newcommand{\lDgDb}{{\mathversion{fordg}%
\resizebox{!}{12pt}{\ATextClr{%
\rotatebox{17}{$\varsigma$}}}}\hspace{-4pt}%
\resizebox{!}{6.5pt}{\ATextClr{D\hspace{-2pt}B}}}}}

\newcommand{\URL}{\resizebox{!}{7pt}{\AcronymText{URL}}}
\newcommand{\CSML}{\resizebox{!}{7pt}{\AcronymText{CSML}}}
\newcommand{\LPF}{\resizebox{!}{7pt}{\AcronymText{LPF}}}
\newcommand{\lLPF}{\resizebox{!}{8.5pt}{\AcronymText{LPF}}}
\newcommand{\lsLPF}{\resizebox{!}{9.5pt}{\AcronymText{LPF}}}

\newcommand{\AI}{\resizebox{!}{7.5pt}{\AcronymText{AI}}}
\newcommand{\lAI}{\resizebox{!}{8pt}{\AcronymText{AI}}}

\newcommand{\Jupyter}{\resizebox{!}{7pt}{\AcronymText{Jupyter}}}
\newcommand{\Python}{\resizebox{!}{7pt}{\AcronymText{Python}}}
\newcommand{\IDN}{\resizebox{!}{7pt}{\AcronymText{IDN}}}
\newcommand{\JPG}{\resizebox{!}{7pt}{\AcronymText{JPG}}}
\newcommand{\JPEG}{\resizebox{!}{7pt}{\AcronymText{JPEG}}}
\newcommand{\PNG}{\resizebox{!}{7pt}{\AcronymText{PNG}}}
\newcommand{\TIFF}{\resizebox{!}{7pt}{\AcronymText{TIFF}}}
\newcommand{\REPL}{\resizebox{!}{7pt}{\AcronymText{REPL}}}

\newcommand{\Pandore}{\resizebox{!}{7pt}{\AcronymText{Pandore}}}

\newcommand{\MIFlowCyt}{\resizebox{!}{7pt}{\AcronymText{MIFlowCyt}}}
\newcommand{\GatingML}{\resizebox{!}{7pt}{\AcronymText{Gating-ML}}}
\newcommand{\flowCL}{\resizebox{!}{7pt}{\AcronymText{flowCL}}}


\makeatletter

\newcommand*\getX[1]{\expandafter\getX@i#1\@nil}

\newcommand*\getY[1]{\expandafter\getY@i#1\@nil}
\def\getX@i#1,#2\@nil{#1}
\def\getY@i#1,#2\@nil{#2}
\makeatother
	
\newcommand{\rectann}[9]{%
\path [draw=#1,draw opacity=#2,line width=#3, fill=#4, fill opacity = #5, even odd rule] %
(#6) rectangle(\getX{#6}+#7,\getY{#6}+#8)
({\getX{#6}+((#7-(#7*#9))/2)},{\getY{#6}+((#8-(#8*#9))/2)}) rectangle %
({\getX{#6}+((#7-(#7*#9))/2)+#7*#9},{\getY{#6}+((#8-(#8*#9))/2)+#8*#9});}


\definecolor{pfcolor}{RGB}{94, 54, 73}

\newcommand{\EPF}{\resizebox{!}{7pt}{\AcronymText{ETS{\color{pfcolor}pf}}}}
\newcommand{\lEPF}{\resizebox{!}{8.5pt}{\AcronymText{ETS{\color{pfcolor}pf}}}}
\newcommand{\lsEPF}{\resizebox{!}{9.5pt}{\AcronymText{ETS{\color{pfcolor}pf}}}}


\newcommand{\XPDF}{\resizebox{!}{7pt}{\AcronymText{XPDF}}}

\newcommand{\GRE}{\resizebox{!}{7pt}{\AcronymText{GRE}}}
\newcommand{\CAS}{\resizebox{!}{7pt}{\AcronymText{CAS}}}

\newcommand{\lMOSAIC}{%
\resizebox{!}{8pt}{\ATexttclr{M}}%
\resizebox{!}{6pt}{\ATexttclr{OSAIC}}}

\newcommand{\llMOSAIC}{%
\resizebox{!}{8.5pt}{\ATexttclr{MOSAIC}}}

\newcommand{\ConceptsDB}{\resizebox{!}{7pt}{\ATexttclr{ConceptsDB}}}


\newcommand{\XML}{\resizebox{!}{7pt}{\AcronymText{XML}}}

\newcommand{\lXML}{\resizebox{!}{8pt}{\AcronymText{XML}}}

\newcommand{\RDF}{\resizebox{!}{7pt}{\AcronymText{RDF}}}
\newcommand{\DOM}{\resizebox{!}{7pt}{\AcronymText{DOM}}}

\newcommand{\Java}{\resizebox{!}{7pt}{\AcronymText{Java}}}


\newcommand{\ParaView}{\resizebox{!}{7pt}{\AcronymText{ParaView}}}
\newcommand{\Octave}{\resizebox{!}{7pt}{\AcronymText{Octave}}}
\newcommand{\ROOT}{\resizebox{!}{7pt}{\AcronymText{ROOT}}}
\newcommand{\CERN}{\resizebox{!}{7pt}{\AcronymText{CERN}}}
\newcommand{\MQFour}{\resizebox{!}{7pt}{\AcronymText{MQ4}}}
\newcommand{\VISSION}{\resizebox{!}{7pt}{\AcronymText{VISSION}}}

\newcommand{\sVISSION}{\resizebox{!}{6pt}{\AcronymText{VISSION}}}


\newcommand{\ReproZip}{\resizebox{!}{7pt}{\AcronymText{ReproZip}}}
\newcommand{\BioCoder}{\resizebox{!}{7pt}{\AcronymText{BioCoder}}}

\newcommand{\Covid}{\resizebox{!}{7pt}{\AcronymText{Covid-19}}}


\newcommand{\HMCL}{{\resizebox{!}{7.5pt}{\ATexttclr{HMCL}}}}
\newcommand{\DSPIN}{{\resizebox{!}{7.5pt}{\ATexttclr{D-SPIN}}}}

\newcommand{\Pandore}{{\resizebox{!}{7.5pt}{\ATexttclr{Pandore}}}}

\newcommand{\lDSPIN}{{\resizebox{!}{8pt}{\ATexttclr{D-SPIN}}}}
\newcommand{\lsDSPIN}{{\resizebox{!}{8.5pt}{\ATexttclr{D-SPIN}}}}

\newcommand{\CLang}{\resizebox{!}{7pt}{\AcronymText{C}}}

\newcommand{\HNaN}{\resizebox{!}{7pt}{\AcronymText{HN%
\textsc{a}N}}}

\newcommand{\JSON}{\resizebox{!}{7pt}{\AcronymText{JSON}}}
\newcommand{\UV}{\resizebox{!}{7pt}{\AcronymText{UV}}}


\newcommand{\PET}{\resizebox{!}{7pt}{\AcronymText{PET}}}
\newcommand{\MRI}{\resizebox{!}{7pt}{\AcronymText{MRI}}}


\newcommand{\MeshLab}{\resizebox{!}{7pt}{\AcronymText{MeshLab}}}
\newcommand{\IQmol}{\resizebox{!}{7pt}{\AcronymText{IQmol}}}

\newcommand{\SGML}{\resizebox{!}{7pt}{\AcronymText{SGML}}}

\newcommand{\WhiteDB}{\resizebox{!}{7pt}{\AcronymText{\makebox{WhiteDB}}}}

\newcommand{\CrossRef}{\resizebox{!}{7pt}{\AcronymText{CrossRef}}}

\newcommand{\ASCII}{\resizebox{!}{7pt}{\AcronymText{ASCII}}}

\newcommand{\GUI}{\resizebox{!}{7pt}{\AcronymText{GUI}}}
\newcommand{\UI}{\resizebox{!}{7pt}{\AcronymText{UI}}}


\newcommand{\URI}{\resizebox{!}{7pt}{\AcronymText{URI}}}
\newcommand{\DTD}{\resizebox{!}{7pt}{\AcronymText{DTD}}}

\newcommand{\API}{\resizebox{!}{7pt}{\AcronymText{API}}}

\newcommand{\JATS}{\resizebox{!}{7pt}{\AcronymText{JATS}}}


\newcommand{\SDI}{\resizebox{!}{7pt}{\AcronymText{SDI}}}
\newcommand{\SDIV}{\resizebox{!}{7pt}{\AcronymText{SDIV}}}

\definecolor{atColor}{RGB}{50, 22, 40}
\newcommand{\ATextClr}[1]{\textcolor{atColor}{\textbf{#1}}}

\newcommand{\DgDb}{\makebox{\raisebox{-3pt}{\resizebox{!}{11pt}{\ATextClr{%
\rotatebox{17}{$\varsigma$}}}}\hspace{-4pt}%
\resizebox{!}{6.5pt}{\ATextClr{D\hspace{-2pt}B}}}}

\newcommand{\lDgDb}{\makebox{\raisebox{-3pt}{%
\resizebox{!}{12pt}{\ATextClr{%
\rotatebox{17}{$\varsigma$}}}}\hspace{-4pt}%
\resizebox{!}{6.5pt}{\ATextClr{D\hspace{-2pt}B}}}}


\newcommand{\IDE}{\resizebox{!}{7pt}{\AcronymText{IDE}}}

\newcommand{\OWL}{\resizebox{!}{7pt}{\AcronymText{OWL}}}

\newcommand{\Kaggle}{\resizebox{!}{7pt}{\AcronymText{Kaggle}}}


\newcommand{\ViSion}{\resizebox{!}{7pt}{\AcronymText{ViSion}}}

\newcommand{\CWL}{\resizebox{!}{7pt}{\AcronymText{CWL}}}

\newcommand{\ThreeD}{\resizebox{!}{7pt}{\AcronymText{3D}}}
\newcommand{\TwoD}{\resizebox{!}{7pt}{\AcronymText{2D}}}

\newcommand{\medInria}{\resizebox{!}{7pt}{\AcronymText{medInria}}}
\newcommand{\ThreeDimViewer}{\resizebox{!}{7pt}{\AcronymText{3DimViewer}}}

\newcommand{\FAIR}{\resizebox{!}{7pt}{\AcronymText{FAIR}}}

\newcommand{\MIAPEGI}{\resizebox{!}{7pt}{\AcronymText{MIAPE-GI}}}
\newcommand{\MIAPE}{\resizebox{!}{7pt}{\AcronymText{MIAPE}}}



\newcommand{\QNetworkManager}{\resizebox{!}{7pt}{\AcronymText{QNetworkManager}}}
\newcommand{\QTextDocument}{\resizebox{!}{7pt}{\AcronymText{QTextDocument}}}
\newcommand{\QWebEngineView}{\resizebox{!}{7pt}{\AcronymText{QWebEngineView}}}
\newcommand{\HTTP}{\resizebox{!}{7pt}{\AcronymText{HTTP}}}


\newcommand{\lAcronymTextNC}[2]{{\fontfamily{fvs}\selectfont {\Large{#1}}{\large{#2}}}}

\newcommand{\AcronymTextNC}[1]{{\fontfamily{fvs}\selectfont {\large #1}}}


\colorlet{orr}{orange!60!red}

\newcommand{\textscc}[1]{{\color{orr!35!black}{{%
						\fontfamily{Cabin-TLF}\fontseries{b}\selectfont{\textsc{\scriptsize{#1}}}}}}}


\newcommand{\textsccserif}[1]{{\color{orr!35!black}{{%
				\scriptsize{\textbf{#1}}}}}}


\newcommand{\iXPDF}{\resizebox{!}{7pt}{\textsccserif{%
\textit{XPDF}}}}

\newcommand{\iEPF}{\resizebox{!}{7pt}{\textsccserif{%
\textit{ETSpf}}}}

\newcommand{\iSDI}{\resizebox{!}{7pt}{\textsccserif{%
\textit{SDI}}}}

\newcommand{\iHTXN}{\resizebox{!}{7pt}{\textsccserif{%
\textit{HTXN}}}}


\newcommand{\AcronymText}[1]{{\textscc{#1}}}

\newcommand{\AcronymTextser}[1]{{\textsccserif{#1}}}


\newcommand{\mAcronymText}[1]{{\textscc{\normalsize{#1}}}}

\newcommand{\FASTA}{{\resizebox{!}{7pt}{\AcronymText{FASTA}}}}
\newcommand{\SRA}{{\resizebox{!}{7pt}{\AcronymText{SRA}}}}
\newcommand{\DNA}{{\resizebox{!}{7pt}{\AcronymText{DNA}}}}
\newcommand{\MAP}{{\resizebox{!}{7pt}{\AcronymText{MAP}}}}
\newcommand{\EPS}{{\resizebox{!}{7pt}{\AcronymText{EPS}}}}
\newcommand{\CSV}{{\resizebox{!}{7pt}{\AcronymText{CSV}}}}
\newcommand{\PDB}{{\resizebox{!}{7pt}{\AcronymText{PDB}}}}

\newcommand{\WebGL}{{\resizebox{!}{7pt}{\AcronymText{WebGL}}}}
\newcommand{\Docker}{{\resizebox{!}{7pt}{\AcronymText{Docker}}}}


\newcommand{\OBO}{{\resizebox{!}{7pt}{\AcronymText{OBO}}}}

\newcommand{\XOCS}{{\resizebox{!}{7pt}{\AcronymText{XOCS}}}}

\newcommand{\ChemXML}{{\resizebox{!}{7pt}{\AcronymText{ChemXML}}}}

\newcommand{\TeXMECS}{\resizebox{!}{7pt}{\AcronymText{TeXMECS}}}

% pmml  arff  openannotation

\newcommand{\PMML}{\resizebox{!}{7pt}{\AcronymText{PMML}}}
\newcommand{\ARFF}{\resizebox{!}{7pt}{\AcronymText{ARFF}}}
\newcommand{\IeXML}{\resizebox{!}{7pt}{\AcronymText{IeXML}}}

\newcommand{\SeCo}{\resizebox{!}{7pt}{\AcronymText{SeCo}}}


\newcommand{\HDFFive}{\resizebox{!}{7pt}{\AcronymText{HDF5}}}

\newcommand{\NGML}{\resizebox{!}{7pt}{\ATexttclr{NGML}}}


\newcommand{\Cpp}{\resizebox{!}{7pt}{\AcronymText{C++}}}

%\newcommand{\\WhiteDB{}}{\resizebox{!}{7pt}{\AcronymText{\WhiteDB{}}}}

\colorlet{drp}{darkRed!70!purple}

%\newcommand{\MOSAIC}{{\color{drp}{\AcronymTextNC{\scriptsize{MOSAIC}}}}}

\newcommand{\MOSAIC}{\resizebox{!}{7pt}{\ATexttclr{MOSAIC}}}


\newcommand{\mMOSAIC}{{\color{drp}{\ATexttclr{\normalsize{MOSAIC}}}}}

\newcommand{\lmMOSAIC}{\ATexttclr{\Large{M}\normalsize{OSAIC}}}

\newcommand{\MOSAICVM}{\mMOSAIC-\AcronymTextNC{VM}}
\newcommand{\MOSAICSR}{\resizebox{!}{7pt}{%
\mMOSAIC-\AcronymTextNC{\textcolor{drp!25!black}{SR}}}}

\newcommand{\MOSAICSD}{\resizebox{!}{7pt}{%
\mMOSAIC-\AcronymTextNC{\textcolor{drp!25!black}{SD}}}}

\newcommand{\MOSAICGIS}{\resizebox{!}{7pt}{%
\mMOSAIC-\AcronymTextNC{\textcolor{drp!25!black}{GIS}}}}


\newcommand{\lMOSAICSR}{\resizebox{!}{7.5pt}{%
\lmMOSAIC-\AcronymTextNC{\textcolor{drp!25!black}{SR}}}}

\newcommand{\lMOSAICSD}{\resizebox{!}{7pt}{%
\lmMOSAIC-\AcronymTextNC{\textcolor{drp!25!black}{SD}}}}


\newcommand{\lsMOSAICSR}{\resizebox{!}{9.5pt}{%
\lMOSAIC-}\resizebox{!}{9pt}{\AcronymTextNC{\textcolor{drp!45!black}{SR}}}}

%\newcommand{\lMOSAICSR}{\resizebox{!}{7pt}{%
%\mMOSAIC-\AcronymTextNC{SR}}}


\newcommand{\sMOSAICVM}{\resizebox{!}{7pt}{\MOSAICVM}}
\newcommand{\sMOSAIC}{\resizebox{!}{7pt}{\MOSAIC}}

\newcommand{\LDOM}{\resizebox{!}{7pt}{\AcronymText{LDOM}}}
\newcommand{\Cnineteen}{\resizebox{!}{7pt}{\AcronymText{CORD-19}}}

\newcommand{\lCnineteen}{\resizebox{!}{7.5pt}{\AcronymText{CORD-19}}}


\newcommand{\MOL}{\resizebox{!}{7pt}{\AcronymText{MOL}}}

\newcommand{\ACL}{\resizebox{!}{7pt}{\AcronymText{ACL}}}

\newcommand{\LXCR}{\resizebox{!}{7pt}{\AcronymText{LXCR}}}
\newcommand{\lLXCR}{\resizebox{!}{8.5pt}{\AcronymText{LXCR}}}
\newcommand{\lsLXCR}{\resizebox{!}{9.5pt}{\AcronymText{LXCR}}}

%\newcommand{\lMOSAIC}{{\color{drp}{\lAcronymTextNC{M}{OSAIC}}}}
\newcommand{\lfMOSAIC}{\resizebox{!}{9pt}{{\color{drp}{\lAcronymTextNC{M}{OSAIC}}}}}

\newcommand{\Mosaic}{\resizebox{!}{7pt}{\MOSAIC}}
\newcommand{\MosaicPortal}{{\color{drp}{\AcronymTextNC{MOSAIC Portal}}}}

\newcommand{\RnD}{\resizebox{!}{7pt}{\AcronymText{R\&D}}}

\newcommand{\MIBBI}{\resizebox{!}{7pt}{\AcronymText{MIBBI}}}

\newcommand{\JVM}{\resizebox{!}{7pt}{\AcronymText{JVM}}}
\newcommand{\ECL}{\resizebox{!}{7pt}{\AcronymText{ECL}}}

\newcommand{\ChaiScript}{\resizebox{!}{7pt}{\AcronymText{ChaiScript}}}

\newcommand{\TCP}{\resizebox{!}{7pt}{\AcronymText{TCP}}}

\newcommand{\lQt}{\resizebox{!}{8.5pt}{\AcronymText{Qt}}}
\newcommand{\QtCpp}{\resizebox{!}{8.5pt}{\AcronymText{Qt/C++}}}
\newcommand{\Qt}{\resizebox{!}{7pt}{\AcronymText{Qt}}}

\newcommand{\QtSQL}{\resizebox{!}{7pt}{\AcronymText{QtSQL}}}

\newcommand{\HTML}{\resizebox{!}{7pt}{\AcronymText{HTML}}}
\newcommand{\PDF}{\resizebox{!}{7pt}{\AcronymText{PDF}}}

\newcommand{\R}{\resizebox{!}{7pt}{\AcronymText{R}}}
\newcommand{\SciXML}{\resizebox{!}{7pt}{\AcronymText{SciXML}}}

\newcommand{\MPF}{\resizebox{!}{7pt}{\ATexttclr{MPF}}}


\newcommand{\lGRE}{\resizebox{!}{7.5pt}{\AcronymText{GRE}}}

\newcommand{\p}[1]{

\vspace{.9em}#1}

\newcommand{\q}[1]{{\fontfamily{qcr}\selectfont ``}#1{\fontfamily{qcr}\selectfont ''}} 

%\newcommand{\deconum}[1]{{\textcircled{#1}}}

\renewcommand{\thesection}{\protect\hspace{-1.5em}}
%\renewcommand{\thesection}{\protect\mbox{\deconum{\Roman{section}}}}
\renewcommand{\thesubsection}{\protect\hspace{-1em}}

\newcommand{\llMOSAIC}{\mbox{{\LARGE MOSAIC}}}
%\newcommand{\lfMOSAIC}{\mbox{M\small{OSAIC}}}

\newcommand{\llMosaic}{\llMOSAIC}
\newcommand{\lMosaic}{\lMOSAIC}
\newcommand{\lfMosaic}{\lfMOSAIC}

%\newcommand{\dsC}{}

\newcommand{\textds}[1]{{\fontfamily{lmdh}\selectfont{%
\raisebox{-1pt}{#1}}}}

\newcommand{\ltextds}[1]{{\fontfamily{lmdh}\fontsize{12}{11}\selectfont{%
\raisebox{-1pt}{#1}}}}

\newcommand{\dsC}{{\textds{ds}{\fontfamily{qhv}\selectfont \raisebox{-1pt}{C}}}}
\newcommand{\ldsC}{{\textds{ds}{\fontfamily{qhv}\selectfont \raisebox{-1pt}{C}}}}

\newcommand{\MdsX}{\resizebox{!}{9pt}{\ATexttclr{\raisebox{-1pt}{{\fontfamily{lmdh}\selectfont M}}\fontfamily{qhv}\selectfont dsX}}}
\newcommand{\lsMdsX}{\resizebox{!}{10.5pt}{\ATexttclr{\raisebox{-1pt}{{\fontfamily{lmdh}\selectfont M}}\fontfamily{qhv}\selectfont dsX}}}
\newcommand{\lMdsX}{\resizebox{!}{9.5pt}{\ATexttclr{\raisebox{-1pt}{{\fontfamily{lmdh}\selectfont M}}\fontfamily{qhv}\selectfont dsX}}}


\newcommand{\llWC}{\mbox{{\LARGE WhiteCharmDB}}}

\newcommand{\llwh}{\mbox{{\LARGE White}}}
\newcommand{\llch}{\mbox{{\LARGE CharmDB}}}

\usepackage{enumitem}
%\usepackage{listings}

\colorlet{dsl}{purple!20!brown}
\colorlet{dslr}{dsl!50!blue}

\setlist[description]{%
  topsep=11pt,
  labelsep=22pt, leftmargin=10pt,
  itemsep=9pt,               % space between items
  %font={\bfseries\sffamily}, % set the label font
  font=\normalfont\bfseries\color{dslr!50!black}, % if colour is needed
}

\setlist[enumerate]{%
  topsep=3pt,               % space before start / after end of list
  itemsep=-2pt,               % space between items
  font={\bfseries\sffamily}, % set the label font
%  font={\bfseries\sffamily\color{red}}, % if colour is needed
}

%\usepackage{tcolorbox}

\newcommand{\slead}[1]{%
\noindent{\raisebox{2pt}{\relscale{1.15}{{{%
\fcolorbox{logoCyan!50!black}{logoGreen!5}{#1}
}}}}}\hspace{.5em}}


\let\OldLaTeX\LaTeX

\renewcommand{\LaTeX}{\resizebox{!}{7pt}{\color{orr!35!black}{\OldLaTeX}}}

\let\OldTeX\TeX

\renewcommand{\TeX}{\resizebox{!}{7pt}{\color{orr!35!black}{\OldTeX}}}


\newcommand{\LargeLaTeX}{\resizebox{!}{8.5pt}{\color{orr!35!black}{\OldLaTeX}}}

\setlength\parindent{0pt}
%\setlength\parindent{24pt}
%
\usepackage{titlesec}
\usepackage[hyphens]{url}

\newsavebox{\tempbox}
\sbox{\tempbox}{\raisebox{-4.5pt}%
{\parbox{3cm}{\textit{\textbf{{\large ``Channel \hspace{1em} Abstractions"}}}}}}

%\newsavebox{\sstwolinebox}
%\sbox{\sstwolinebox}{\raisebox{-5pt}%
%{\parbox{7.4cm}{\textit{\textbf{{\large %
%Range-Bounded Types, Value Constructors, and Addressability}}}}}}
%\newcommand{\sstwoline}{\usebox{\sstwolinebox}}

%\definecolor{midRed}{rgb}{.5,.0,.25}
%\definecolor{logoRed}{rgb}{.3,0,0}



\newenvironment{displayquotexxx}{%
	%\begin{fcolorbox}{yellow!20!gray}{red!5}
	\begin{flushright}
	\begin{tcolorbox}[
		colback=gray!4,
		colframe=red!20!gray,
		width=.98\linewidth,
		boxrule=1pt,leftrule=3pt,
		rightrule=1pt,toprule=.5pt,arc=0pt,auto outer arc
		]
\begin{minipage}{21.5em}}{%
\end{minipage}\end{tcolorbox}
	\end{flushright}
}

\usepackage{changepage}

\newenvironment{dquote}{%
	%\begin{fcolorbox}{yellow!20!gray}{red!5}
	\begin{flushright}
	\vspace{1em}\begin{tcolorbox}[
		breakable, parbox=false, colback=gray!4,
		colframe=mRed!20!gray,
		width=.995\linewidth,
		boxrule=1pt,leftrule=3pt,
		rightrule=1pt,toprule=.5pt,arc=0pt,auto outer arc
		]
\begin{adjustwidth}{1pt}{-3pt}
}{%
\end{adjustwidth}
\end{tcolorbox}\vspace{1em}
	\end{flushright}
}



\newcounter{subs}[section]

\newcommand{\subss}[2]{	
\phantomsection \label{#1}
\addcontentsline{toc}{subsection}{#1}
\stepcounter{subs}
\refstepcounter{subsection}
\vspace*{3.25em}

\noindent{ {\Large\textbf\thesection.}{\large\textbf\thesubs}} {
 {{\Large\textbf #2}}
}
\vspace*{.35em} 
}

\newcommand{\subsstlx}[2]{}

\newcommand{\subsstl}[2]{	
\phantomsection \label{#1}
\addcontentsline{toc}{subsection}{#1}
\stepcounter{subs}
\refstepcounter{subsection}
\vspace*{3.25em}

\noindent{ \raisebox{-1pt}{\Large\textbf\thesection.}{\large\textbf\thesubs}} {
\raisebox{-1pt}{{\Large\textbf #2}}
}
\vspace*{.35em} 
}

\newcommand{\subsstly}[2]{	
\phantomsection \label{#1}
\addcontentsline{toc}{subsection}{#1}
\stepcounter{subs}
\refstepcounter{subsection}
\vspace*{3.25em}

\noindent{ \raisebox{-1pt}{\Large\textbf\thesection.}{\large\textbf\thesubs}} {
\raisebox{-1pt}{{\Large\textbf #2}}
}
\vspace*{.35em} 
}


\let\Osubsection\subsection

\renewcommand{\subsection}[1]{\subss{#1}{#1}}

\newcommand{\subsectionalt}[2]{\subss{#1}{#2}}


\newsavebox{\twolinebox}

\newcommand{\stwoline}[1]{%
\sbox{\twolinebox}{\raisebox{-3pt}%
{\parbox{7.4cm}{\linespread{1.25}\selectfont\raggedleft{\textbf{{\large #1}}}}}}}

\newcommand{\twoline}{\usebox{\twolinebox}}

\newcommand{\subsectiontwoline}[1]{\stwoline{#1}\subsstl{#1}{\twoline}

\vspace{1em}}

\newcommand{\spsubsectiontwolinerepl}[2]{\stwoline{#1}\subsstl{#2}{\twoline}}
\newcommand{\subsectiontwolinerepl}[2]{\stwoline{#1}\subsstl{#2}{\twoline}

\vspace{1em}}

\newcommand{\subsectiontwolinealt}[2]{\stwoline{#1}\subsstl{#2}{\twoline}

\vspace{1em}}


\titlespacing{\subsection}{0pt}{20pt}{20pt}
\titlespacing{\section}{0pt}{35pt}{15pt}
\titlespacing{\subsubsection}{0pt}{20pt}{5pt}


\newcommand{\spsubsection}[1]{%
\subss{#1}{#1}
\vspace{1em}
}
\newcommand{\spsubsectiontwoline}[1]{%
\subsectiontwoline{#1}
\vspace{1em}
}



\usepackage{csquotes}

\usepackage{booktabs}

\usepackage{amssymb}

\usepackage{titlesec}

\usepackage{transparent}

\usepackage{setspace}

\usepackage{graphicx}

\newcommand{\sectionGraphics}{

\vspace{2em}\centerline{\includegraphics[scale=0.125]{logo-deco.png}}\vspace{-1em}}

\usepackage[flushmargin]{footmisc}

\setlength{\parindent}{30pt}

\usepackage[letterpaper, left=.47in,right=.47in,top=.825in,bottom=.825in]{geometry}

\let\xbibitem\bibitem
\renewcommand{\bibitem}[2]{\vspace{6.5pt}\xbibitem{#1}{#2}}

\newenvironment{frquote}{%
%\begin{fcolorbox}{yellow!20!gray}{red!5}
%\begin{flushright}

\begin{tcolorbox}[
	colback=white,
	colframe=white,
	width=.97\linewidth,
	arc=0mm, auto outer arc
	]
\begin{scriptsize}
\begin{minipage}{61em}
\begin{flushright}
\begin{minipage}{63em}}{%
\end{minipage}
\end{flushright}
\end{minipage}
\end{scriptsize}
\end{tcolorbox}
%\end{flushright}
}

%\colorlet{codegr}{black!80!blue}

\setlength{\columnsep}{8mm}

\usepackage{etoolbox}

\AtBeginEnvironment{thebibliography}{\linespread{1}\selectfont}

\usepackage{mathptmx}

\titleformat*{\subsection}{\small\bfseries}

\usepackage{wasysym}
\usepackage{textcomp}
\usepackage{amssymb}

\usepackage{microtype}

\DeclareMathAlphabet{\mathcal}{OMS}{cmsy}{m}{n}

\let\OldI\i

\newcommand{\secvspace}{\vspace{-0.3em}}
\newcommand{\asecvspace}{\vspace{-0.05em}}

\newcommand{\mdash}{---}
\newcommand{\q}[1]{``#1"}
\newcommand{\sq}[1]{`#1'}
\renewcommand{\i}[1]{\textit{#1}}

\newcommand{\nl}{

}

\newcommand{\T}[1]{\raisebox{-2pt}{\ensuremath{\mathcal{T}}}\textit{\tiny #1}}

\newcommand{\biburl}[1]{ {\fontfamily{gar}\selectfont{\textcolor[rgb]{.2,.6,0}%
{\scriptsize {\url{#1}}}}}}

%\newcommand{\biburl}[1]{ {\fontfamily{gar}\selectfont{\textcolor[rgb]%%{.2,.6,0}%
%{\scriptsize \textls*[-70]{\burl{#1}}}}}}

\newcommand{\TSupT}{\ensuremath{{\T2}\makebox[4pt][r]{\raisebox{5pt}{{\scalebox{.6}{\T1}}}}}}

\let\OldFootnoteSize\footnotesize
\renewcommand{\footnotesize}{\scriptsize}

\newcommand{\emigres}{\'emigr\'es}

\newcommand{\SExpressions}{S-Expressions}
\newcommand{\SExpression}{S-Expression}

\newcommand{\cq}[1]{{\fontfamily{gar}\selectfont ``}#1{\fontfamily{gar}\selectfont "}}

\newcommand{\Retore}{Retor\'e}
\newcommand{\Aurelie}{Aur\'elie}
\newcommand{\Descles}{D\'escles}

\newcommand{\SG}{\ensuremath{\mathbb{SG}}}
\newcommand{\tLa}{\ensuremath{\mathbb{A}_\lambda}}

\newcommand{\lt}{\ensuremath{<}}
\newcommand{\ancestorLt}{\ensuremath{\lll}}
\newcommand{\aALTb}{\ensuremath{a {\lll} b}}
\newcommand{\aLTb}{\ensuremath{a {\lt} b}}
\newcommand{\aLTcLtb}{\ensuremath{a {\lt} b {\lt} c}}
\newcommand{\aDLTb}{\ensuremath{a {\lessdot} b}}

\newcommand{\aMath}{\ensuremath{a}}
\newcommand{\bMath}{\ensuremath{b}}
\newcommand{\cMath}{\ensuremath{c}}

\newcommand{\ala}{\`a la}

\newcommand{\br}{

}


\newcommand{\CppEleven}{{\Cpp}11}

\usepackage{graphicx}

\usepackage[breakable]{tcolorbox}

\newsavebox\lstbox



\newif\iffootnote
\let\Footnote\footnote
\renewcommand\footnote[1]{\begingroup\footnotetrue\Footnote{#1}\endgroup}

\colorlet{orr}{orange!60!red}

\newcommand{\textscc}[1]{{\color{orr!35!black}{{%
						\fontfamily{Cabin-TLF}\fontseries{b}\selectfont{\textsc{#1}}}}}}

\newcommand{\textscnc}[1]{{%
						\fontfamily{Cabin-TLF}\fontseries{b}\selectfont{\textsc{#1}}}}

\newcommand{\AcronymTextNC}[1]{{\iffootnote\begin{footnotesize}{\textscnc{#1}}\end{footnotesize}%
\else\begin{small}{\textscnc{#1}}\end{small}\fi}}

\newcommand{\AcronymText}[1]{{\iffootnote\begin{scriptsize}{\textscc{#1}}\end{scriptsize}%
\else\begin{OldFootnoteSize}{\textscc{#1}}\end{OldFootnoteSize}\fi}}

\newcommand{\smAcronymText}[1]{\begin{footnotesize}\textscc{#1}\end{footnotesize}}

\newcommand{\librets}{\AcronymText{librets}}
\newcommand{\FHIR}{\AcronymText{FHIR}}
\newcommand{\DICOM}{\AcronymText{DICOM}}

\newcommand{\IoT}{\AcronymText{IoT}}


\newcommand{\TCP}{\AcronymText{TCP}}
\newcommand{\GPS}{\AcronymText{GPS}}
\newcommand{\CppTwenty}{\AcronymText{C++20}}

\newcommand{\AcronymTextInitialCap}[1]{{\iffootnote\begin{scriptsize}{\textsc{#1}}\end{scriptsize}%
		\else\begin{normalsize}{\textsc{#1}}\end{normalsize}\fi}}


\newcommand{\rel}[1]{\raisebox{0.25pt}{%
	{\iffootnote\begin{footnotesize}{\textsc{#1}}\end{footnotesize}%
		\else\begin{OldFootnoteSize}{\textsc{#1}}\end{OldFootnoteSize}\fi}}}

\newcommand{\vs}{

\vspace*{0.2em}}


\newcommand{\HTML}{\AcronymText{HTML}}

\newcommand{\UX}{\AcronymText{UX}}

\newcommand{\Csh}{\AcronymText{C\#}}

\newcommand{\PHP}{\AcronymText{PHP}}
\newcommand{\IDE}{\AcronymText{IDE}}

\newcommand{\AND}{\AcronymText{AND}}
\newcommand{\DX}{\AcronymText{DX}}

\newcommand{\RelaeGraph}{\AcronymText{RelaeGraph}}
\newcommand{\PDF}{\AcronymText{PDF}}

\newcommand{\NDP}{\AcronymText{NDP}}
\newcommand{\UML}{\AcronymText{UML}}
\newcommand{\GIT}{\AcronymText{GIT}}

\newcommand{\UI}{\AcronymText{UI}}
\newcommand{\QML}{\AcronymText{QML}}

\newcommand{\SPARQL}{\AcronymText{SPARQL}}
\newcommand{\OWL}{\AcronymText{OWL}}
\newcommand{\JSON}{\AcronymText{JSON}}

\newcommand{\const}{\AcronymText{const}}
\newcommand{\GeCODE}{\AcronymText{GeCODE}}

\newcommand{\Python}{\AcronymText{Python}}


\newcommand{\GUI}{\AcronymText{GUI}}
\newcommand{\BNF}{\AcronymText{BNF}}


\newcommand{\NLP}{\AcronymText{NLP}}

\newcommand{\Francois}{Fran\c{}cois}


\newcommand{\Gardenfors}{G\"ardenfors}


\newcommand{\Haskell}{\AcronymText{Haskell}}

\newcommand{\Cpp}{\AcronymText{C++}}
\newcommand{\RDF}{\AcronymText{RDF}}
\newcommand{\Java}{\AcronymText{Java}}
\newcommand{\IT}{\AcronymText{IT}}
\newcommand{\AI}{\AcronymText{AI}}

\newcommand{\SCA}{\AcronymText{SCA}}

\newcommand{\Lisp}{\AcronymText{Lisp}}

\newcommand{\Scheme}{\AcronymText{Scheme}}

\newcommand{\ThreeD}{\AcronymText{3D}}

\definecolor{DarkRed}{rgb}{.2,.0,.1}

\colorlet{orrr}{orange!40!red}
\colorlet{orrbl}{orrr!85!blue}
\colorlet{orrb}{orrbl!80!DarkRed}

\newcommand{\lclc}[1]{{\color{orrb}{#1}}}

\newcommand{\lcl}[2]{{\resizebox{!}{#1}{\color{orrb}{#2}}}}
\newcommand{\ty}{{\lcl{8pt}{\ensuremath{\mathfrak{t}}}}}

\newcommand{\caltypeT}{\ensuremath{\ty}}
\newcommand{\calS}{\ensuremath{\mathcal{S}}}

\newcommand{\typeTp}{\lclc{\ensuremath{\ty'}}}
\newcommand{\typeTpp}{\lclc{\ensuremath{\ty''}}}


\newcommand{\gOpTransferOneOneF}{\codeText{g:}\codeTextr{return$_1$}{\opTransfer}%
\codeText{f:}\codeTextr{lambda$_1$}}
 
\newcommand{\fDotOfg}{\codeText{f$\circ$g}}
\newcommand{\fDotOfGX}{\codeText{(f.g)(x)}}
\newcommand{\fOfGx}{\codeText{f(g(x))}}

\newcommand{\fgx}{\codeText{f(g(x))}}

\newcommand{\funG}{\codeText{g}}
\newcommand{\funF}{\codeText{f}} 

\newcommand{\tyOne}{\codeText{${\ty}_1$}}
\newcommand{\tyTwo}{\codeText{${\ty}_2$}}

\newcommand{\tyOneTotyTwo}{\tyOne \codeText{$\rightarrow$} \tyTwo}

\newcommand{\chK}{\codeText{$\mathcal{K}$}}

\newcommand{\tOneTimesTTwo}{\codeText{${\ty}_1$ $\times$ ${\ty}_2$}}
\newcommand{\tOneTimesTOne}{\codeText{${\ty}_1$ $\times$ ${\ty}_1$}}
\newcommand{\tyOneTimesTyTwo}{\codeText{${\ty}_1$ $\times$ ${\ty}_2$}}

\newcommand{\tOneTimesTTwoToTOneOntoTTwo}{\codeText{${\ty}_1$ $\times$ ${\ty}_2$ %
$\Longrightarrow$ ${\ty}_1$ $\rightarrow$ ${\ty}_2$}}

\newcommand{\unitTy}{\codeText{Unit}} 
\newcommand{\unitVal}{\codeText{unit}} 

\newcommand{\unitTyToty}{\unitTy{}{ }%
\codeText{$\rightarrow$ {\ty}}}

\newcommand{\tyOneToTyTwo}{\codeText{${\ty}_1$ $\rightarrow$ ${\ty}_2$}}

\newcommand{\tyE}{\lclc{$E$}}
\newcommand{\tyTotyE}{\codeText{{\ty} $\rightarrow$ $E$}}
\newcommand{\tyToTyE}{\codeText{{\ty} $\rightarrow$ $E$}}
 
\newcommand{\tyValues}{{\ty}-values}


	
\newcommand{\Tnoindex}{\raisebox{-2pt}{\ensuremath\ty}}

\newcommand{\typeAbove}{%
\raisebox{-1pt}{\rotatebox{90}{\begin{tiny}$\diagdown$\makebox[1pt][c]{$\diagup$}\end{tiny}}}}

\newcommand{\typeT}{\ensuremath{type\raisebox{.5pt}{\makebox[3pt][c]{-}}\ty}}
\newcommand{\TValues}{{\ensuremath\ty}-values} 

\newcommand{\tOnetoTwotoThree}{\codeText{${\ty}_1$\smsp%
$\rightarrow$\smsp${\ty}_2$\smsp$\rightarrow$\smsp${\ty}_3$}}
 
\newcommand{\tOnetoTwoTOThree}{\codeText{${\ty}_1$\smsp%
$\rightarrow$\smsp\smsp(${\ty}_2$\smsp$\rightarrow$\smsp${\ty}_3$)}} 


\definecolor{BaseColor}{HTML}{8533FF}

\colorlet{ftcfore}{BaseColor!60!cyan}
\colorlet{ftcback}{BaseColor!40!cyan}

\newcommand{\ftc}[3]{
\vspace*{6mm}
\begin{tcolorbox}
[float=t, colframe=ftcfore!20!white,boxrule=0.5pt,arc=22pt,enhanced,
title={{\color{black}{\protect{#1}}}},label={#2}
toprule=0pt,bottomrule=1pt,
drop fuzzy shadow northeast={darkRed!40!ftcback},
      boxsep=3pt]\hspace{3em}\parbox{0.8\textwidth}\protect{#3}
\end{tcolorbox}      
\vspace*{-4mm}
}

\newcommand{\tc}[2]{
\vspace*{6mm}
\begin{tcolorbox}
[#1 colframe=darkRed!70!BaseColor,boxrule=0.5pt,arc=22pt,enhanced,
toprule=0pt,bottomrule=1pt,
drop fuzzy shadow northeast={darkRed!70!purple},
      boxsep=3pt]\hspace{3em}\parbox{0.8\textwidth}\protect{#2}
\end{tcolorbox}      
\vspace*{-4mm}
}

\tcbuselibrary{skins}
\usetikzlibrary{calc}
\usetikzlibrary{shadows}
\pgfdeclarelayer{background}
\pgfdeclarelayer{foreground}
\pgfsetlayers{background,main,foreground}



\usetikzlibrary{fit}
\usepackage{caption}

\renewcommand{\figurename}{Diagram}

\newcommand{\emblink}[2]{\href{#1}{#2}}

\newcommand{\tmphs}{\hypersetup{linkbordercolor=orange!50!red,linkcolor=black}}
\newcommand{\tmphscol}{\hypersetup{linkbordercolor=gray!40,linkcolor=black}}

\usepackage[object=vectorian]{pgfornament} 

\newcommand{\decoline}{\vspace{-4em}}

\newcommand{\decolinex}{%

\vspace{-2em}
{\color{darkRed!60!cyan}\noindent\hfil{\EnglischeLinie}\hfil}
\vspace{-2.25em}}


\newcommand{\whdecoline}{%

\vspace{-2em}
{\color{white}\noindent\hfil{\whEnglischeLinie}\hfil}
\vspace{-2.25em}}

\newcommand{\sectionline}[1]{%
  \noindent
  \begin{center}
  {\color{#1}
    \resizebox{0.5\linewidth}{1ex}
    {{%
    {\begin{tikzpicture}
    \node  (C) at (0,0) {};
    \node (D) at (9,0) {};
    \path (C) to [ornament=84] (D);
    \end{tikzpicture}}}}}%
    \end{center}
  }
  
\newcommand{\EnglischeLinie}{
\sectionline{darkRed!60!cyan}
}

\newcommand{\whEnglischeLinie}{
\sectionline{white}
}

\newcommand{\thinsectionline}[1]{%
	\noindent
	\begin{center}
		{\color{#1}
			\resizebox{.2\linewidth}{1.5ex}
			{{%
					{\begin{tikzpicture}
						\node  (C) at (0,0) {};
						\node (D) at (9,0) {};
						\path (C) to [ornament=84] (D);
						\end{tikzpicture}}}}}%
	\end{center}
}

\newcommand{\thindecolinex}{\vspace*{-.15em}\thinsectionline{blGreen!10!orange}\vspace*{-.45em}}
\newcommand{\thindecolineadjx}[2]{\vspace*{#1}\thinsectionline{blGreen!10!orange}\vspace*{#2}}


\newcommand{\thindecoline}{\vspace*{-.15em}\thinsectionline{black!70}\vspace*{-.45em}}
\newcommand{\thindecolineadj}[2]{\vspace*{#1}\thinsectionline{black!70}\vspace*{#2}}


\newlength{\mplength}
\setlength{\mplength}{1.05\linewidth}

\newlength{\hslength}
\setlength{\hslength}{-1.45cm}

\newcommand{\spinc}[1]{\input{#1}}

\newsavebox{\tcsb}


\newcommand{\spinctc}[3]{\begin{lrbox}{\tcsb}\protect\input{#1}\end{lrbox} %\input{#1}
\begin{figure}
\tc{}{\protect\usebox\tcsb} 
\captionof{figure}{#2}
\label{#3}
\end{figure}
}

\usetikzlibrary{backgrounds} 
\usetikzlibrary{arrows}
\tikzset{>=triangle 90}

\newcommand{\gFunB}{\ensuremath{\mathfrak{g}}}
\newcommand{\fFunB}{\ensuremath{\mathfrak{f}}}

%\newcommand{\lstinlinebstyle}[1]{\Colorbox{blue!10}{%
%		{\fontfamily{Cabin-TLF}\fontseries{b}\fontsize{9pt}{16pt}\selectfont{{\textls[200]{#1}}}}}}

\newcommand{\lstinlinebstyle}[1]{\textls[200]{#1}}


\newenvironment{tikzgrid}{%
\begin{tikzpicture}[shorten >=1pt,node distance=2cm,on grid,auto,
framed,background rectangle/.style={double,ultra thick,draw=darkRed, bottom color=cyan!20, 
	top color=black!20, rounded corners}]%
}{\end{tikzpicture}}
 
\usepackage[font=small,labelfont=bf]{caption}

\definecolor{BaseColor}{HTML}{8533FF}

\colorlet{ftcfore}{BaseColor!60!cyan}
\colorlet{ftcback}{BaseColor!40!cyan}


\newcommand{\tcl}[1]{\begin{lrbox}{\lstbox}\begin{minipage}{\mplength}
\input{#1}
\end{minipage}\end{lrbox}
\tc{}{\hspace*{\hslength}\usebox\lstbox}
}

\newcommand{\itcl}[1]{
\input{#1}
}

\newcommand{\itclfig}[2]{
\begin{figure}\input{#1}
\label{#2}
\end{figure}
}

\newcommand{\ftcl}[2]{\begin{lrbox}{\lstbox}\begin{minipage}{\mplength}
\input{#2}
\end{minipage}\end{lrbox}
\tc{float=#1,}{\hspace*{\hslength}\usebox\lstbox}
}

\newcommand{\vsftcl}[2]{\begin{lrbox}{\lstbox}\begin{minipage}{\mplength}
\vsinput{#2}
\end{minipage}\end{lrbox}
\tc{float=#1,}{\hspace*{\hslength}\usebox\lstbox}
}

\newcommand{\vstcl}[1]{\begin{lrbox}{\lstbox}\begin{minipage}{\mplength}
%\vsinput{#1}
\end{minipage}\end{lrbox}  
%\tc{}{\hspace*{\hslength}\usebox\lstbox}
}


\definecolor{blGreen}{rgb}{.2,.7,.3}
\definecolor{darkRed}{rgb}{.2,.0,.1}

\definecolor{darkBlGreen}{rgb}{.1,.3,.2}

\definecolor{oldBlColor}{rgb}{.2,.7,.3}

\definecolor{blColor}{rgb}{.1,.3,.2}

\definecolor{elColor}{rgb}{.2,.1,0}
\definecolor{flColor}{rgb}{0.7,0.3,0.3}

\definecolor{logoOrange}{RGB}{108, 18, 30}
\definecolor{logoGreen}{RGB}{85, 153, 89}
\definecolor{logoPurple}{RGB}{200, 208, 30}

\definecolor{logoBlue}{RGB}{4, 2, 25}
\definecolor{logoPeach}{RGB}{255, 159, 102}
\definecolor{logoCyan}{RGB}{66, 206, 244}
\definecolor{logoRed}{rgb}{.3,0,0}

\definecolor{mRed}{rgb}{.5,.0,.25}



\definecolor{Bkg}{RGB}{250,245,252}
\newcommand{\leader}[2]{\hspace{#1}\colorbox{Bkg}{#2}}


\newcommand{\saying}[1]{\vspace{2ex}\noindent{%%
				\leader{2em}{\begin{minipage}{.93\textwidth}{\footnotesize #1}\end{minipage}}\vspace{2ex}}}

\newcommand{\sayingsrc}[1]{\vspace{0ex}\\\hspace{2pt} --- {\footnotesize #1}}

\renewcommand{\labelitemi}{{\raisebox{4pt}{{\footnotesize{$\bullet$}}}}}

\newcommand{\itemmark}{\raisebox{-4pt}{\rotatebox{90}{{\Large $\bracevert$}}}}

\usepackage{tikz}
\usetikzlibrary{positioning}
\usetikzlibrary{shapes,snakes}

\newcommand{\visavis}{vis-\`a-vis}

\newcommand{\tinyurl}[1]{{\raisebox{2pt}{{\scriptsize \url{#1}}}}} 

\usepackage[colorlinks=true]{hyperref}

\usepackage{eso-pic}

\colorlet{urlclr}{red!40!magenta!50!orange}

\hypersetup{
 urlcolor = urlclr,
 urlbordercolor = cyan!60!black,
 linkcolor = red!30!black,
 citecolor = orange!30!black,
 citebordercolor = yellow!30!black,
} 



\newcommand{\p}[1]{
	
	\vspace{.65em}
	#1	
}

\let\OldEnumerate\enumerate
  

\let\OldSection\section

\renewcommand\section[1]{
	\vspace{12pt}
	
	\scalebox{1.3}{\colorbox{logoPurple!50}{\hspace{1em}}}
					
	{\protect\transparent{0.5}{\colorbox{logoPeach}{%
			\begin{minipage}{\linewidth}
				    \vspace{.5em}
				
	\protect\transparent{1}{\OldSection{#1%
     }}
	    \end{minipage}}} }
    
    \vspace{-5em}
    	{\protect\transparent{0.3}{\colorbox{logoPeach}{%
    		\begin{minipage}{\linewidth}
    			\hspace{\linewidth}
    	\end{minipage}}} } 
    
    \vspace{5em}
    
	\vspace{-6pt}
}


\newcommand\tabltxt[1]{\footnotesize{#1}}
\renewcommand\section[1]{\OldSection{#1}}

\titleformat*{\subsection}{\Large\bfseries}

%\let\OldSubsection\subsection
%\renewcommand\subsection[1]{

%	\vspace{12pt}
	
%		\OldSubsection{% 	
%			\hspace{-2.75em}
%			\protect\raisebox{-5pt}{%
%			\colorbox{logoCyan!50}{\hspace{2.1em}}}%
%			\hspace{-5pt}{\protect\transparent{0.3}{\colorbox{logoBlue!80}{\protect\transparent{1}{%
%						   \protect\raisebox{1pt}{\textit{{\large #1}}} }}}}
%		}
%	\vspace{-6pt}
%}


\let\OLDthebibliography\thebibliography
\renewcommand\thebibliography[1]{
\let\section\OldSection
\setlength{\leftmargin}{-4pt}
\vspace{.1em}
\OLDthebibliography{#1}
\vspace{.7em}
\OldFootnoteSize 
\setlength{\parskip}{0pt}
\setlength{\itemsep}{1pt plus 0ex}
\raggedright
}

\makeatletter
\def\@biblabel#1{\hspace{-6pt}#1}
\makeatother

\newcommand{\bibtitle}[1]{{\small \textit{#1}}}
\newcommand{\intitle}[1]{{\hspace{3pt}\textls*[-80]{\texttt{\textit{#1}}}}\hspace{-1pt}}

\renewcommand{\i}[1]{\textit{#1}}

\usepackage{enumitem}
\setlist{leftmargin=1mm}



\newcommand{\itemtitle}[1]{{\color{green!10!red!40!black} \textls*[-80]{\texttt{#1}}}}

\newcommand{\acronymText}[1]{\AcronymText{#1}}
\newcommand{\acronymTextInitialCap}[1]{\AcronymTextInitialCap{#1}}
\newcommand{\acronymTextSS}[1]{\textsc{#1}}



\newcommand{\MIT}{\AcronymText{MIT}}

\newcommand{\nulltt}{\AcronymText{\texttt{null}}}

\newcommand{\boostUnits}{\codeText{\texttt{boost::units}}}

\newcommand{\Maybe}{\codeText{Maybe}}
\newcommand{\bind}{\codeText{bind}}
\newcommand{\return}{\codeText{return}}

\newcommand{\elseif}{\codeText{else if}}

\newcommand{\IDL}{\AcronymText{IDL}}

\newcommand{\CoAP}{\AcronymText{CoAP}}
\newcommand{\MQTT}{\AcronymText{MQTT}}

\usepackage{tabularx}

\newcommand{\UDPipe}{\AcronymText{UDPipe}}
\newcommand{\mWSI}{\AcronymText{mWSI}}

\newcommand{\smRGBa}{\smAcronymText{RGBa}}
\newcommand{\smHSV}{\smAcronymText{HSV}}
\newcommand{\smRGB}{\smAcronymText{RGB}}

\newcommand{\OS}{\AcronymText{OS}}

\newcommand{\DTD}{\AcronymText{DTD}}

\newcommand{\XML}{\AcronymText{XML}}
\newcommand{\API}{\AcronymText{API}}
\newcommand{\IR}{\AcronymText{IR}}
\newcommand{\Clang}{\AcronymText{Clang}}

\newcommand{\HCI}{\AcronymText{HCI}}
\newcommand{\HTTP}{\AcronymText{HTTP}}
\newcommand{\CSS}{\AcronymText{CSS}}

\newcommand{\STL}{\AcronymText{STL}}

\newcommand{\Qt}{\AcronymText{Qt}}
\newcommand{\NL}{\AcronymText{NL}}

\newcommand{\vsinput}[1]{\vspace{1em}\input{#1}}

\newcommand{\qi}[1]{\q{\textit{#1}}}

\newcommand{\CSharp}{\AcronymText{C\#}}
%\newcommand{\IoT}{\AcronymText{IoT}}

\newcommand{\CoNLL}{\AcronymText{CoNLL}}
\newcommand{\CoNLLU}{\AcronymText{CoNLL-U}}

\usetikzlibrary{decorations.pathmorphing}

\usepackage{listings}

\renewcommand{\lstlistingname}{Sample}

%%%%%%%%
\newcommand{\OO}{\AcronymText{OO}}
\newcommand{\SQL}{\AcronymText{SQL}}
\newcommand{\QtCpp}{\AcronymText{Qt/C++}}
\newcommand{\JavaScript}{\AcronymText{JavaScript}}

\newcommand{\ECL}{\AcronymText{ECL}}

\newcommand{\ovn}[1]{\color{yellow}{{\textbf{#1}}}}

\newcommand{\sovn}[1]{\color{yellow}{{\textbf{#1}}}}

\newcommand{\dovn}[3]{\draw[draw=blue,fill=DarkRed] (#1,#2) circle[radius=3mm];
\node at (#1,#2){\ovn{#3}}
}

\newcommand{\sdovn}[3]{\draw[draw=blue,fill=DarkRed] (#1,#2) circle[radius=1.4mm];
\node (char) at (#1,#2) {\sovn{#3}}
}

\newcommand{\sdovnz}[3]{\raisebox{-1pt}{{\draw[draw=blue,fill=DarkRed] (#1,#2) circle[radius=1.4mm];
\node (char) at (#1,#2) {\sovn{#3}}
}}}

\newcommand*\circledx[1]{\tikz[baseline=(char.base), inner sep=0]{
            \sdovn{0}{0}{#1};}}

\newcommand{\circled}[1]{\raisebox{-.5pt}{\circledx{#1}}}

\newcommand{\circledup}[1]{\circledx{#1}}

\newcommand{\circledd}[1]{\raisebox{-2pt}{\circledx{#1}}}
			
\newcommand{\OneOverlay}{\circled{1}}
\newcommand{\TwoOverlay}{\circled{2}}
\newcommand{\ThreeOverlay}{\circled{3}}
\newcommand{\FourOverlay}{\circled{4}}
\newcommand{\FiveOverlay}{\circled{5}}
\newcommand{\SixOverlay}{\circled{6}}

\newcommand{\OneOverlayu}{\circledup{1}}
\newcommand{\TwoOverlayu}{\circledup{2}}
\newcommand{\ThreeOverlayu}{\circledup{3}}
\newcommand{\FourOverlayu}{\circledup{4}}
\newcommand{\FiveOverlayu}{\circledup{5}}
\newcommand{\SixOverlayu}{\circledup{6}}

\newcommand{\FourOverlayd}{\circledd{4}}

\newcommand{\true}{\codeText{true}}

\colorlet{codegr}{black!80!blue}

\newcommand{\FnDoc}{\codeText{Fn\_Doc}}

\newcommand{\KCMEnv}{\codeText{KCM\_Env}}

\newcommand{\envv}{\codeText{envv}}

\newcommand{\kprom}{\codeText{kcm\_promote\_type\_binding\_to%
\_statement\_via\_type\_de\-fault\_literal\_\_let}}
\newcommand{\kcmde}{\codeText{kcm\_direct\_eval}}

\newcommand{\dbleq}{\codeText{==}}
\newcommand{\Sdbleq}{\codeText{$\backslash$==}}
\newcommand{\Seq}{\codeText{$\backslash $=}}

\usepackage[utf8]{inputenc}
 
\usepackage{listings}
\usepackage{color}
 
\definecolor{codegreen}{rgb}{0,0.6,0}
\definecolor{codegray}{rgb}{0.5,0.5,0.5}
\definecolor{codepurple}{rgb}{0.58,0,0.82}
\definecolor{backcolour}{rgb}{0.95,0.95,0.92}
 
\usetikzlibrary{positioning}

\usepackage{listings}
\lstset{%
  frame            = tb,    % draw frame at top and bottom of code block
  tabsize          = 1,     % tab space width
  numbers          = left,  % display line numbers on the left
  framesep         = 3pt,   % expand outward
  framerule        = 0.4pt, % expand outward 
  commentstyle     = \color{Green},      % comment color
  keywordstyle     = \color{blue},       % keyword color
  stringstyle      = \color{DarkRed},    % string color
  backgroundcolor  = \color{gray!20}, % backgroundcolor color
  showstringspaces = false,              % do not mark spaces in strings
  upquote          = true
}

\newcommand{\yeqfxz}{\codeText{y=f(x,z)}}

\newcommand{\try}{\codeText{try}}
\renewcommand{\DH}{\AcronymText{DH}}
\newcommand{\CH}{\AcronymText{CH}}

%\let{\oldeta}{\eta}
%\renewcommand{\eta}{\ensuremath{eta}}


\newcommand{\catch}{\codeText{catch}}


\newcommand{\vecrgbt}{\AcronymText{RGBT}}
\newcommand{\vecrgb}{\AcronymText{RGB}}

\newcommand{\TTL}{\AcronymText{TTL}}

\newcommand{\URL}{\AcronymText{URL}}
\newcommand{\RRF}{\AcronymText{RRF}}
\newcommand{\SPO}{\angla{\AcronymText{Subject}, %
\AcronymText{Predicate}, \AcronymText{Object}}}

\newcommand{\Turtle}{\AcronymText{Turtle}}

\newcommand{\CLang}{\AcronymText{C}}

 
\newcommand{\intxeqninety}{\codeText{int x = 90}}

\newcommand{\tyFile}{\codeText{file}}
\newcommand{\idrisText}[1]{\codeText{#1}}
\newcommand{\openFn}{\codeText{open}}

\newcommand{\ML}{\acronymText{ML}}

\newcommand{\largeRDF}{RDF}

\newcommand\Small{\fontsize{8}{10}\selectfont}
\newcommand*\LSTfont{%
  \Small\ttfamily\SetTracking{encoding=*}{-50}\lsstyle}

	
\newcommand{\NathLivedTTL}{\lstset{ basicstyle=\LSTfont, columns=fullflexible, xleftmargin=5mm, framexleftmargin=5mm, numbers=left, stepnumber=1, breaklines=true, breakatwhitespace=false, numberstyle=\footnotesize, numbersep=5pt, 
tabsize=2, frame=lines, captionpos=b, caption={Turtle Formats}}
    %\lstinputlisting{NathLivedTTL.n3}
    }		

\newcommand{\TyS}{\codeTextr{$\mathbb{T}$}}

\newcommand{\TXLTyS}{$\codeTextr{\mathfrak{L}_\codeTextr{\mathbb{T}}}$}

\newcommand{\TXLTySChi}{$\codeTextr{\mathfrak{L}_\codeTextr{\mathbb{T}}\chiussr}$}


\newcommand{\ageFF}{\codeTexto{${\lceil}45{\rceil}$}}
\newcommand{\Nath}{\codeTexto{${\lceil}$Nathaniel${\rceil}$}}

\newcommand{\nodeNOne}{$N_1$}
\newcommand{\nodeNTwo}{$N_2$}

\newcommand{\NThree}{\AcronymText{N3}}

\newcommand{\ceila}[1]{${\lceil}$#1${\rceil}$}
\newcommand{\angla}[1]{${\langle}$#1${\rangle}$}

\newcommand{\NathFF}{\codeTexto{\angla{\ceila{Nathaniel}, \ceila{46}}}}
\newcommand{\NathFFBD}{\codeTexto{\angla{\ceila{Nathaniel}, \ceila{46}, %
\ceila{Brooklyn}, \ceila{Democrat}}}}

\newcommand{\BrookDem}{\codeTexto{\angla{\ceila{Brooklyn}, \ceila{Democrat}}}}

\newcommand{\nameAge}{\codeTexto{${\langle}$\codeText{name}, \codeText{age}${\rangle}$}}

\newcommand{\suigeneris}{\i{sui generis}}
 
\newcommand{\struct}{\codeText{struct}}
\newcommand{\CStructs}{\codeText{struct}s}
\newcommand{\CStructsArrays}{\codeText{struct}s/arrays}

\newcommand{\CStrucst}{C\codeText{struct}}

\newcommand{\throw}{\codeText{throw}}
\newcommand{\exception}{\codeText{exception}}

\newcommand{\float}{\codeText{float}}

\newcommand{\chiuss}{\raisebox{-1pt}{$^\chiu$}}
\newcommand{\chiussr}{\raisebox{-1pt}{$^\chiur$}}

\newcommand{\chiur}{\codeTextr{\ensuremath{\chi}}}
\newcommand{\chiu}{\codeText{\ensuremath{\chi}}}

\newcommand{\TySChi}{\TyS\chiussr}


\newcommand{\lCh}{\chsnt{lam}}
\newcommand{\rCh}{\chsnt{ret}}
\newcommand{\xCh}{\chsnt{exc}}

\newcommand{\lrCh}[1]{\chsnt{lam, ret, #1}}

\newcommand{\lrChwow}[1]{\chsnt{lam!, ret!, #1}}

\newcommand{\lsrxCh}{\chsnt{lam, sig, ret, exc}}

\newcommand{\lrChblank}{\chsnt{lam, ret}}

\newcommand{\rChSize}{\chsntsz{ret}}
\newcommand{\chChSize}{\chsntsz{ch}}

\newcommand{\rNoMixx}{\rCh{\chcolor{$\nshortparallel$}}\xCh}
\newcommand{\rOneOrx}{\rCh{\chcolor{$\asymp$\xCh}}}

\newcommand{\rChSizeleOne}{\codeText{\rChSize{} $\leq$ 1}}
\newcommand{\rChSizeLe}{\codeText{rChSizele}}
\newcommand{\Ch}{\chsnt{ch}}

\newcommand{\lrxSimple}{\lrChwow{exc?}}

\newcommand{\lrxDetailed}{\makebox{%
\chsnt{lam!*},\chnt{ret!\%},\chnt{exc?\%}}}

\newcommand{\excl}{\chcolor{\ensuremath{\asymp}}}

\newcommand{\lrxTotal}{\makebox{%
\lrCh{exc} {\colonblg} %
\chnt{lam!*},\chnt{ret!\%},\chnt{exc?\%}%
{ }{\colonblg}{\colonblg} \chnt{ret}{\excl}\chnt{exc}}}

\newcommand{\lr}{\makebox{%
\lrChblank{}{ }{\colonblg} \chnt{lam!*},\chnt{ret!\%}}}

\newcommand{\lsrx}{\makebox{\makebox{\lsrxCh}%
{ }{\colonblg} \chnt{lam!*},\chnt{sig?\%},\chnt{ret!\%},\chnt{exc?\%}}}
\newcommand{\sCh}{\chsnt{sig}}


\newcommand{\lrx}{\makebox{%
\lrCh{exc}}}

\let\OldLambda\lambda

\renewcommand{\lambda}{\codeText{lambda}}
\renewcommand{\return}{\codeText{return}}

\newcommand{\lambdaPLUSreturn}{{\lambda}{\codeTextr{+}}{\return}}


\newcommand{\lxr}{\codeText{lxr}}

\newcommand{\fnote}[1]{\codeText{#1}}

\newcommand{\codeinclude}{\codeText{\#include}}
\newcommand{\codeconstruct}{\codeText{construct}}

\newcommand{\USH}{\acronymText{USH}}

\newcommand{\fxy}{\codeText{\makebox{f(x, y)}}}
\newcommand{\fSym}{\codeText{f}}

\newcommand{\fofx}{\codeText{\makebox{f(x)}}}
\newcommand{\fofy}{\codeText{\makebox{f(y)}}}

\newcommand{\fFuns}{\codeText{f}}
\newcommand{\hFun}{\codeText{h}}
\newcommand{\hfx}{\codeText{h(f(x))}}
\newcommand{\fxdoth}{\codeText{f(x).h()}}

\newcommand{\typesH}{\codeText{types.h}}

\newcommand{\QDataStream}{\codeText{QDataStream}}

\newcommand{\CATSet}{$\mathbb{S}$}

\newcommand{\DSL}{\AcronymText{DSL}}
\newcommand{\TXL}{\AcronymText{TXL}}
\newcommand{\CXL}{\AcronymText{CXL}}

\newcommand{\IntZToOH}{\codeText{int\ranged{0,100}}}

\newcommand{\OH}{\codeText{100}}

\newcommand{\codebreak}{\codeText{break}}
\newcommand{\codecontinue}{\codeText{continue}}

\newcommand{\TVOneToVTwo}{\codeText{T\ranged{$V_1$,$V_2$}}}

\newcommand{\RangeLTEVal}{\codeText{ranged\_lte}}
\newcommand{\RangeLTEOHxeqOHone}{\codeText{ranged\_lte{$<$100$>$} x = 101}}
\newcommand{\RangeLTEValV}{\codeText{\codeText{ranged\_lte{$<$V$>$}}}}

\newcommand{\MIpair}{\codeText{mi\_pair}} 
\newcommand{\fmipair}{\codeText{int f(mi\_pair pr)}}

\newcommand{\fileTy}{\codeText{file}} 
\newcommand{\pairOfLists}{\codeText{pair$<$list$<$...$>>$}}


\newcommand{\VOne}{\codeText{$V_1$}}
\newcommand{\VTwo}{\codeText{$V_2$}}

\newcommand{\TType}{\codeText{T}}

\newcommand{\crVOverRTwo}{\codeTextr{$\sqrt{3}\sqrt[3]{V}$}}

\newcommand{\volSphCube}{\codeTextr{$\frac{4}{3}\sqrt{27}\pi V$}}

\newcommand{\rRad}{\codeTextr{$R$}}
\newcommand{\vVol}{\codeTextr{$V$}}
\newcommand{\piSym}{\codeTextr{$\pi$}}


\newcommand{\TMyList}{\codeText{MyList$<$T$>$}}

\newcommand{\MyList}{\codeText{MyList}}

\newcommand{\templateTMyList}{\codeText{template$<$T$>$MyList}}
\newcommand{\MyListInt}{\codeText{MyList$<$int$>$}}

\newcommand{\listsize}{\codeText{list.size()}}
\newcommand{\sizelist}{\codeText{size(list)}}

\newcommand{\SLE}{\AcronymText{SLE}}

\newcommand{\tCat}{\TyS}
\newcommand{\DBH}{\AcronymText{DBH}}

\newcommand{\cCat}{$\mathbb{C}$}
\newcommand{\eOneToeTwo}{$e_1 \rightarrow e_2$}
\newcommand{\eOne}{$e_1$}
\newcommand{\eTwo}{$e_2$}

\newcommand{\VolSphere}{\codeTextr{$\frac{{4\pi R^3 }}{3}$}}

\newcommand{\livedinup}{\raisebox{6pt}{\colorbox{blGreen!20}{lived-in}}}

\newcommand{\rdfrestup}{\raisebox{4pt}{\colorbox{blGreen!20}{rdf:rest}}}


\newcommand{\CArray}{\AcronymText{C} array}
\newcommand{\CStruct}{\AcronymText{C} \codeText{struct}}

\newcommand{\WAV}{\AcronymText{WAV}}
\newcommand{\NPY}{\AcronymText{NPY}}



\newcommand{\Chi}{$\codeTextr{\chi}$}

\newcommand{\NTrips}{\AcronymText{N-Triples}}
\newcommand{\NTrip}{\AcronymText{N-Triple}}

\renewcommand{\GeCODE}{\AcronymText{GeCODE}}
\newcommand{\ITK}{\AcronymText{ITK}}

\newcommand{\SSI}{\AcronymText{SSI}}
\newcommand{\GSI}{\AcronymText{GSI}}

\newcommand{\ISBN}{\AcronymText{ISBN}}
\newcommand{\GIS}{\AcronymText{GIS}}


\newcommand{\CodeMinted}[1]{{\color{codegr}{{%
\fontfamily{lmss}\fontseries{b}\selectfont{#1}}}}}

\newcommand{\CodeMintedo}[1]{{\color{orange!40!black}{{%
\fontfamily{lmss}\fontseries{b}\selectfont{#1}}}}}

\newcommand{\CodeMintedr}[1]{{\color{red!40!black}{{%
\fontfamily{lmss}\fontseries{b}\selectfont{#1}}}}}

\newcommand{\codeText}[1]{\CodeMinted{#1}}
\newcommand{\codeTexto}[1]{\CodeMintedo{#1}}
\newcommand{\codeTextr}[1]{\CodeMintedr{#1}}

\newcommand{\FactPP}{\AcronymText{Fact++}}

\newcommand{\chname}[1]{\AcronymTextNC{\textbf{#1}}}

\newcommand{\chsym}{\raisebox{3pt}{\rotatebox{-45}{$\Arrowvert$}}\hspace{-5pt}%
\raisebox{1pt}{\rotatebox{-45}{\tiny{$\gg$}}}\hspace{-1pt}}  %\searrow

\newcommand{\chsymt}{\raisebox{3pt}{\rotatebox{-45}{$\Arrowvert$}}\hspace{-5pt}%
\raisebox{1pt}{\rotatebox{-45}{\tiny{$\gg$}}}\hspace{-1pt}}  %\searrow \Rightarrow

\newcommand{\chsnt}[1]{{\color{blGreen!40!black}{\chsymt\chname{#1}}}}
\newcommand{\chnt}[1]{{\color{blGreen!40!black}{\chname{#1}}}}

\newcommand{\colonblg}{{\color{blGreen!40!black}{:}}}

\newcommand{\chcolor}[1]{{\color{blGreen!40!black}{#1}}}

\newcommand{\chsntsz}[1]{{\color{blGreen!40!black}{\makebox{\#\chsymt\chname{#1}}}}}


\newcommand{\chsn}[1]{\chsym\chname{#1.}}

\newcommand{\lXY}{\codeTexto{${\OldLambda}x.{\OldLambda}y$}}


\newcommand{\CHlXY}{\chcolor{\chsn{lam}$xy$}}



\newcommand{\Tvar}{\codeText{T}}
\newcommand{\TrRan}{\codeText{T\ranged{r}}}
\newcommand{\rRan}{\codeText{\ranged{r}}}

\newcommand{\jFunction}{\ensuremath{j}-function}

\newcommand{\xVal}{\codeText{x}}
\newcommand{\xeqyplusz}{\codeText{x $=$ y $+$ z}}

\newcommand{\fOfG}{\codeText{$f{\circ}g$}}
\newcommand{\Ofop}{\codeText{$\circ$}}

\newcommand{\inc}{\codeText{inc}}

\newcommand{\zeroNum}{\codeText{$0$}}


\newcommand{\enumleft}{\codeText{$\langle$}}
\newcommand{\enumright}{\codeText{$\rangle$}}

\newcommand{\closed}{\codeText{closed}}
\newcommand{\open}{\codeText{open}}
\newcommand{\nonexistent}{\codeText{nonexistent}}

\newcommand{\codeDH}{\codeText{code}-{\DH}}


\newcommand{\fOfg}{\codeText{$f{\circ}g$}}  


\newcommand{\cfFun}{\codeText{Cf}}
\newcommand{\cf}{\codeText{Cf}}
\newcommand{\zToOH}{\codeText{\ranged{0, 100}}}
\newcommand{\C}{\codeText{C}}
\newcommand{\yplusz}{\codeText{y $+$ z}}
\newcommand{\iVal}{\codeText{$i$}}
\renewcommand{\le}{\codeText{$\leq$}}

\renewcommand{\int}{\codeText{int}}

\newcommand{\tTy}{\codeText{T}}

\newcommand{\rrRanOfTVV}{\codeText{ranged$<$T, t1, t2$>$}}

\newcommand{\nVal}{$n$}

\newcommand{\zTon}{\codeText{\ranged{0, n}}} 

\newcommand{\fOneTwoxeq}{\codeText{$f_1(x)=f_2(x)$}}

\newcommand{\fOne}{\codeText{$f_1$}}
\newcommand{\fTwo}{\codeText{$f_2$}}

\newcommand{\xVar}{\codeText{$x$}}

\newcommand{\fx}{\codeText{f(x)}}

\newcommand{\tOne}{\codeText{${\ty}_1$}}
\newcommand{\tTwo}{\codeText{${\ty}_2$}}

\newcommand{\gSym}{\codeText{g}}
\newcommand{\gradeFn}{\codeText{grade}}
\newcommand{\iSym}{\codeText{i}}

\newcommand{\xninetyyonehgrade}{\codeText{x $=$ 90; y $=$ 100; g $=$ grade(x, y)}}

\newcommand{\tOneTotTwo}{\codeText{${\ty}_1 \Rightarrow {\ty}_2$}}

\newcommand{\zeroToOH}{\codeText{\ranged{0,100}}}

\newcommand{\rrsb}[1]{\raisebox{5pt}{#1}}
\newcommand{\ranged}[1]{\codeText{\rrsb{{\tiny{$\lgroup$}}}#1\rrsb{{\tiny{$\rgroup$}}}}}

\newcommand{\ftytwoh}{\codeTextr{{$\big[40-200\big]$}}}

\newcommand{\ZeroToOneHundred}{\codeText{\ranged{0,100}}}

\newcommand{\rRanOfT}{\codeText{ranged$<$T$>$}}


\newcommand{\xSym}{\codeText{x}}
\newcommand{\ySym}{\codeText{y}}
\newcommand{\zSym}{\codeText{z}}
\newcommand{\fFun}{\codeText{f}}

%%% SLE
\newcommand{\addressOf}{\codeText{address-of}}
\newcommand{\fofg}{\codeText{$f{\circ}g$}}

\newcommand{\smsp}{\hspace{2pt}}
 
\newcommand{\pimath}{\ensuremath{\pi}}

\newcommand{\SCO}{\AcronymText{SCO}}
\newcommand{\RnD}{\AcronymText{R{\&}D}}



\newcommand{\lambdaxfgx}{\codeTextr{$\OldLambda{}x.fgx$}} 
\newcommand{\lambdaxfx}{\codeTextr{$\OldLambda{}x.fx$}} 

\newcommand{\gFun}{\codeText{g}}  

\newcommand{\nodex}{\codeText{$x$}}
\newcommand{\nodef}{\codeText{$f$}}

\newcommand{\sqqq}{{\fontfamily{lmtt}\selectfont{"}}}
\newcommand{\operatorqq}{\codeText{operator\sqqq\sqqq}}

\newcommand{\cCar}{{\lcl{7pt}{\ensuremath{\mathfrak{c}}}}}
\newcommand{\cCarOne}{{\lcl{7pt}{\ensuremath{\mathfrak{c_1}}}}}
\newcommand{\cCarTwo}{{\lcl{7pt}{\ensuremath{\mathfrak{c_2}}}}}

\newcommand{\opTransfer}{\codeTextr{$\looparrowright$}}
\newcommand{\carrOne}{\cCarOne}
\newcommand{\carrTwo}{\cCarTwo}

\newcommand{\catchexce}{\codeText{catch(Exception e)}}

\newcommand{\carrOneOpTransferTwo}{\carrOne{}{\opTransfer}\carrTwo}
\newcommand{\carrOneOpTransferTwolambda}{\carrOne{}{\opTransfer}%
{\textsuperscript{\hspace{-.8em}{\lambdach}}}\carrTwo}


\newcommand{\carrOneOpTransferTworeturn}{\carrOne{}{\opTransfer}%
{\textsuperscript{\hspace{-.8em}{\returnch}}}\carrTwo}

\newcommand{\carrTwoOpTransferOnereturn}{\carrTwo{}{\opTransfer}%
{\textsuperscript{\hspace{-.8em}{\returnch}}}\carrOne}


\newcommand{\carOne}{\cCarOne}

\newcommand{\carTwo}{\cCarTwo} 

\newcommand{\carOnetoTwos}{\carOne \codeTextr{$\twoheadrightarrow$} \carTwo}
\newcommand{\carOnetoTwof}{\carOne \codeTextr{$\rightarrowtail$} \carTwo}


\newcommand{\intieqzero}{\codeText{int i = 0}}

\newcommand{\intthrtwo}{\codeText{int32}}

\newcommand{\aeqb}{\codeText{a = b}}

\newcommand{\aceqb}{\codeText{a := b}}

\newcommand{\thisc}{\codeText{this}}

\newcommand{\lambdaCalculus}{$\OldLambda$-Calculus}

\newcommand{\lambdas}{$\OldLambda$s}


\newcommand{\fntoch}{\codeTextr{$f_n$ $\rightarrow$ $\chi$}}

\newcommand{\sigmac}{\codeTextr{sigma}}
\newcommand{\sigmach}{\codeTextr{sigma}} 

\newcommand{\sigmaCalculi}{$\varsigma$-calculi}
\newcommand{\sigmaCalculii}{$\varsigma$-calculi}
\newcommand{\sigmaCalculus}{$\varsigma$-calculus}

\newcommand{\returnc}{\codeTextr{return}}

\newcommand{\lambdach}{\codeTextr{lambda}}

\newcommand{\catchddd}{\codeTextr{catch(...)}}

\newcommand{\sectsym}{\S}

\newcommand{\error}{\codeTextr{exception}}
\newcommand{\errorc}{\codeTextr{exception}}
\newcommand{\errorrc}{\codeTextr{exception}}
\newcommand{\returnrc}{\codeTextr{return}}

\newcommand{\returnch}{\codeTextr{return}}
\newcommand{\fgroundch}{\codeTextr{fground}}


\newcommand{\lambdac}{\codeTextr{lambda}}
\newcommand{\capturec}{\codeTextr{capture}}
\newcommand{\capturech}{\codeTextr{capture}}

\newcommand{\errorch}{\codeTextr{exception}}
\newcommand{\breakch}{\codeTextr{break}}
 
\renewcommand{\break}{\codeText{break}}
\newcommand{\breakct}{\codeText{break}}

\newcommand{\yeqfx}{\codeText{y = f(x)}} 
\newcommand{\frcchoprctocy}{\codeText{frcchoprctocy}} 

\newcommand{\frc}{\codeText{frc}} 
\newcommand{\fFn}{\codeText{fFn}} 

\newcommand{\bindFn}{\codeText{bindFn}} 

\newcommand{\lookupch}{\codeTextr{lookup}} 


\newcommand{\stdFuture}{\codeText{std::future}} 
 
 
\newcommand{\doH}{\codeText{do}}

\newcommand{\ifthenelse}{\codeText{if...then...else}}

\newcommand{\tyListOfInt}{\codeText{list$<$int$>$}}
\newcommand{\tyint}{\codeText{int}}

\let\OldGamma\Gamma

\renewcommand{\Gamma}{\codeTextr{$\OldGamma$}}
\newcommand{\gammaOne}{\codeTextr{$\OldGamma_1$}}
\newcommand{\gamaTwo}{\codeTextr{$\OldGamma_2$}}
\newcommand{\gammaTwo}{\codeTextr{$\OldGamma_2$}}

\newcommand{\pfx}{\codeText{(*f)(x)}}

\newcommand{\oprctoc}{\codeText{oprctoc}}
\newcommand{\oprctoch}{\codeText{oprctoch}}
\newcommand{\choprctoc}{\codeText{choprctoc}}

\newcommand{\carOnetoTwofOP}{\codeText{carOnetoTwofOP}}

\newcommand{\thisob}{\codeText{this}}
\newcommand{\selfob}{\codeText{self}}

\newcommand{\typetoch}{\codeTextr{%
\raisebox{2pt}{$\leftharpoonup$}\hspace{-.9em}\raisebox{1pt}{$\leftharpoondown$}}}

\newcommand{\nodeT}{\codeTextr{T$_{Node}$}}
\newcommand{\predTY}{\codeTextr{T$_{Pred}$}}
\newcommand{\predTYOp}{\codeTextr{\lambdach. \typetoch T$_{Pred}$ %
\returnch. \typetoch T$_{Pred}$}}

\newcommand{\predTYsig}{\sigmach\codeTextr{\%}\returnch\codeTextr{\%}}

\newcommand{\QFile}{\codeText{QFile}}
  

\newcommand{\xCommaY}{\codeText{$x,y$}}  
\newcommand{\xlty}{\codeText{x $<$ y}}   
\newcommand{\unsignedint}{\codeText{unsigned int}}
 
\newcommand{\subjectPredicateObject}{\codeTexto{Subject Predicate Object}}
\newcommand{\rdfType}{\codeTexto{rdf:type}}

\newcommand{\parentOf}{\codeTexto{parentOf}}
\newcommand{\childOf}{\codeTexto{childOf}}
\newcommand{\inverseOp}{\codeTexto{inverseOp}}

\newcommand{\capabilitych}{\codeTextr{capability}} 
 
\newcommand{\addFun}{\codeText{add}} 

\newcommand{\sortFun}{\codeText{sort}}
\newcommand{\sortfn}{\codeText{sort}} 

\newcommand{\incimpl}{\codeText{int inc(int x)\{return add(x,1)\}}} 
 
\newcommand{\addOne}{$\langle$\codeText{\&add, 1}$\rangle$}

\newcommand{\Cfr}{$\langle$\codeText{\&Cf, r$_1$, r$_2$}$\rangle$}


\newcommand{\this}{\codeText{this}} 
 
\newcommand{\fsVal}{\codeText{fs}} 
\newcommand{\static}{\codeText{static}}
 
\newcommand{\vVal}{\codeText{V}}  
\newcommand{\lte}{\codeText{$\leq$}}
\newcommand{\lteVal}{\codeText{$\leq{ }$V}}

\newcommand{\intXeqOPointFive}{\codeText{int x = $0.5$}} 

\newcommand{\RangeGTVal}{\codeText{ranged\_gt}}
\newcommand{\RangeGTValx}{\codeText{ranged\_gt{$<$x$>$}}}

\newcommand{\cFun}{\codeText{C}}

\newcommand{\st}{\codeTextr{s.t}} 

\newcommand{\prVal}{\codeText{pr}}

\newcommand{\parach}{\codeTextr{para}}

\newcommand{\fxypreqMIpairxy}{\codeText{f(int x, int y, mi\_pair pr = mi\_pair(x, y))}}
\newcommand{\fxypreqMIpairZeroOne}{\codeText{f(x, y, mi\_pair(0, 1))}}
 
 
%\newcommand{\fxypreqMIpairxy}{\codeText{f int x, int y, mi\_pair pr = mi\_pair x, y}}
%\newcommand{\fxypreqMIpairZeroOne}{\codeText{def}}
 
 
 
\newcommand{\yVal}{\codeText{y}}
\newcommand{\xSymbol}{\codeText{x}}
 
\newcommand{\subjectnd}{\codeTexto{SUBJECT}}
\newcommand{\predicatend}{\codeTexto{PREDICATE}}
\newcommand{\objectnd}{\codeTexto{OBJECT}}
 
\newcommand{\doNotation}{\codeText{do}-notation}

\newcommand{\Onef}{\codeText{$f_1$}}

\newcommand{\Twof}{\codeText{$f_2$}}

\newcommand{\CLng}{\AcronymText{C}}

\newcommand{\qOneTwoThree}{\q{\codeText{123}}} 
\newcommand{\oneTwoThree}{\codeText{123}}
\newcommand{\ascii}{\AcronymText{ASCII}}
 
 

\newcommand{\exceptionch}{\codeTextr{exception}}

\newcommand{\chanOne}{\codeTextr{$\chi_1$}} 
\newcommand{\chanTwo}{\codeTextr{$\chi_2$}} 

\newcommand{\chanOneOpTransferTwo}{\chanOne{}{\opTransfer}\chanTwo}

\newcommand{\objfx}{\codeText{obj.f(x)}}
\newcommand{\obj}{\codeText{obj}}

%\newcommand{\chanyOne}{\codeTextr{$\chi_1$}} 
%\newcommand{\changTwo}{\codeTextr{$\chi_2$}} 
%\newcommand{\chanyOneOpTransferTwo}{\chanOne{}{\opTransfer}\chanTwo}


\newcommand{\li}{\codeText{li}}

\newcommand{\tyNoun}{\codeText{Noun}}
\newcommand{\tyProposition}{\codeText{Proposition}}

\newcommand{\unicode}{\AcronymText{UNICODE}}
\newcommand{\AOP}{\AcronymText{AOP}} 

\newcommand{\chPkg}{\codeTextr{chPkg}}
\newcommand{\tCatChK}{\codeTextr{tCatChK}}

\newcommand{\sigmaTrans}{\codeText{sigmaTrans}} 

\newcommand{\chan}{\codeText{chan}}  

\newcommand{\ProcOne}{\ensuremath{\mathcal{P}_1}}


\newcommand{\carr}{\codeText{$\mathfrak{C}$}}
\newcommand{\tyEE}{\codeText{$\mathcal{E}$}}

\newcommand{\zero}{\codeText{0}} 
\newcommand{\szEqSizeFile}{\codeText{sz = size(file)}}

\newcommand{\sizeFn}{\codeText{size}} 

\newcommand{\openFileType}{\codeText{Open\_File}} 

\newcommand{\fileSym}{\codeText{file}} 
\newcommand{\szSym}{\codeText{sz}} 

\newcommand{\fileExistsFn}{\codeText{file\_exists}} 

\newcommand{\eli}{\item}


\newcommand{\yEqFx}{\codeText{y = f(x)}}

\newcommand{\bindFnSig}{\makebox{\codeText{($>>=$) :: m a $\Rightarrow$ %
(a $\rightarrow$ m b) $\rightarrow$ m b}}}

\renewcommand{\bindFn}{\codeText{$>>=$}}

\newcommand{\monadA}{\codeText{m a}}

\newcommand{\litOne}{\codeText{1}}
\newcommand{\litOnePtZ}{\codeText{1.0}}
\newcommand{\litOnePtO}{\codeText{1.1}}

\newcommand{\piOne}{\ensuremath{\pi_1}}
\newcommand{\piTwo}{\ensuremath{\pi_2}}
\newcommand{\piThree}{\ensuremath{\pi_3}}

\newcommand{\arrbl}{\ensuremath{\longrightarrow\hspace{-1.32em}\bullet}\hspace{.82em}}
\newcommand{\arrfo}{\ensuremath{\longrightarrow\hspace{-1.52em}\lhd}\hspace{.72em}}

\newcommand{\mGamma}{\ensuremath{\gamma}}

\newcommand{\pisbl}{\piOne{} \arrbl{} \piTwo{}}
\newcommand{\pisfo}{\piOne{} \arrfo{} $\ulcorner\piTwo{},\piThree{}\lrcorner$}



\newcommand{\OHO}{\codeText{101}}
\newcommand{\litOHO}{\codeText{101}}
\newcommand{\litOH}{\codeText{100}}
\newcommand{\litFive}{\codeText{5}}
\newcommand{\five}{\codeText{5}}

\newcommand{\nth}{{$n$}th}
%\newcommand{\nVal}{{$n$}}

\newcommand{\listval}{\codeText{list}}

\newcommand{\rlstsize}{\codeText{\ranged{0, list.size()-1}}}
\newcommand{\intrlstsize}{\codeText{int\ranged{0, list.size()-1}}}

\newcommand{\fIntG}{\codeText{int f(int g)}}
\newcommand{\fIntI}{\codeText{int f(int i)}}

\newcommand{\retFive}{\codeText{return 5}}

\newcommand{\QString}{\codeText{QString}}
\newcommand{\QObject}{\codeText{QObject}}
\newcommand{\ptrv}{\codeText{void*}}

\newcommand{\aVal}{\codeText{a}}

\newcommand{\switch}{\codeText{switch}}

\newcommand{\ztooh}{\codeText{0-100}}

\newcommand{\mOldLambda}{\ensuremath{\OldLambda}}



\renewenvironment{abstract}
 {\small
  \begin{center}
  \bfseries \abstractname\vspace{-.5em}\vspace{0pt}
  \end{center}
  \list{}{%
    \setlength{\leftmargin}{28mm}
    \setlength{\rightmargin}{\leftmargin}
  }
  \item\relax}
 {\endlist}


\newcommand{\pseudoIndent}{

\vspace{-2pt}\hspace{8pt}}
\newcommand{\pseudoInd}{\hspace{30pt}}

\setlist[itemize]{leftmargin=2mm}
\setlist[enumerate]{leftmargin=5mm}

\raggedbottom

\setcounter{secnumdepth}{3}

%\usepackage[capitalise,noabbrev,nameinlink]{cleveref}
%\Crefname{subsection}{Subsection}{Subsections}
%\crefname{subsection}{Subsection}{Subsections}
%\usepackage[none]{hyphenat}




\newcommand{\lun}[1]{\raisebox{-4pt}{\fontfamily{qcr}\selectfont{%
\LARGE{\textbf{\textcolor{tcolor}{#1}}}}}\vspace{-2pt}}

\newcommand{\inditem}{\itemindent10pt\item}

\usepackage{soul}

\definecolor{hlcolor}{RGB}{114, 54, 203}
\colorlet{hlcol}{hlcolor!35}
\sethlcolor{hlcol}

\makeatletter
\def\SOUL@hlpreamble{%
	\setul{}{3ex}%         !!!change this value!!! default is 2.5ex
	\let\SOUL@stcolor\SOUL@hlcolor
	\SOUL@stpreamble
}
\makeatother

\usepackage{scrextend}
%\vspace*{3em}
\newenvironment{mldescription}{\vspace{1em}%
  \begin{addmargin}[4pt]{1em}
    \setlength{\parindent}{-1em}%
    \newcommand*{\mlitem}[1][]{\vspace{5pt}\par\medskip%
%\colorbox{hlcolor}{\textbf{##1}}\quad}\indent
\hl{ \textbf{##1} }\quad}\indent
}{%
  \end{addmargin}
  \medskip
}

\usepackage{marginnote}

\newcommand{\mnote}[1]{%
\vspace*{-2em}
\reversemarginpar
\raisebox{1em}{\marginnote{\parbox{4em}{%
\begin{mdframed}[innerleftmargin=4pt,
	innerrightmargin=1pt,innertopmargin=1pt,
	linecolor=red!20!cyan,userdefinedwidth=4em,
	topline=false,
	rightline=false]
{{\fontfamily{ppl}\fontsize{12}{0}\selectfont
		\textit{#1}}}
\end{mdframed}}
	}[3em]}}

\newcommand{\mnotel}[1]{%
\vspace*{-2em}
\reversemarginpar
\raisebox{-4em}{\marginnote{\parbox{4em}{%
\begin{mdframed}[innerleftmargin=4pt,
	innerrightmargin=1pt,innertopmargin=1pt,
	linecolor=red!20!cyan,userdefinedwidth=4em,
	topline=false,
	rightline=false]
{{\fontfamily{ppl}\fontsize{12}{0}\selectfont
		\textit{#1}}}
\end{mdframed}}
	}[3em]}}

\newcommand{\mnoteh}[3]{%
	\vspace*{#1}
	\reversemarginpar
	\raisebox{#2}{\marginnote{\parbox{4em}{%
				\begin{mdframed}[innerleftmargin=4pt,
					innerrightmargin=1pt,innertopmargin=1pt,
					linecolor=red!20!cyan,userdefinedwidth=4em,
					topline=false,
					rightline=false]
					{{\fontfamily{ppl}\fontsize{12}{0}\selectfont
							\textit{#3}}}
				\end{mdframed}}
			}[3em]}}


\newcommand{\mnoteb}[1]{%
	\vspace*{1em}
	\reversemarginpar
	\raisebox{1em}{\marginnote{\parbox{4em}{%
				\begin{mdframed}[innerleftmargin=4pt,
					innerrightmargin=1pt,innertopmargin=1pt,
					linecolor=red!20!cyan,userdefinedwidth=4em,
					topline=false,
					rightline=false]
					{{\fontfamily{ppl}\fontsize{12}{0}\selectfont
							\textit{#1}}}
				\end{mdframed}}
			}[3em]}}
	
\usepackage{wrapfig}

\usetikzlibrary{arrows, decorations.markings}
\usetikzlibrary{shapes.arrows}

\newcommand{\curicon}[2]{%
	\node at (#1,#2) [
	draw=black,
	%minimum width=2ex,
	inner sep=.7pt,
	fill=white,
	single arrow,
	single arrow head extend=3pt,
	single arrow head indent=1.5pt,
	single arrow tip angle=45,
	line join=bevel,
	minimum height=4.6mm,
	rotate=115
	] {};
}

\makeatletter
\def\@cite#1#2{[\textbf{#1\if@tempswa , #2\fi}]}
\def\@biblabel#1{[\textbf{#1}]}
\makeatother


%\let\origref\ref
%\renewcommand{\ref}[1]{{\LARGE #1}}

%\def\ref#1{\textbf{\origref{{\LARGE #1}}}}

\setlength{\footnotesep}{0pt}

\renewcommand{\thefootnote}{\textcolor{logoGreen!80!logoBlue}{{\fontfamily{qcr}\fontseries{b}\fontsize{10}{4}\selectfont\arabic{footnote}}}}


\newcommand{\LVee}{{\colorbox{cyan!40!yellow}{\textcolor{red!70!navy}{\textbf{\LARGE$\vee$}}}}}
\newcommand{\LWedge}{{\colorbox{cyan!40!yellow}{\textcolor{red!70!navy}{\textbf{\LARGE$\wedge$}}}}}

\renewcommand{\LVee}{}
\renewcommand{\LWedge}{}


\urlstyle{same}

\usepackage[preserveurlmacro]{breakurl}

\newcommand{\bhref}[1]{\href{#1}{\burl{#1}}}

%\setmainfont{QTChanceryType}

\begin{document}

\setlength{\skip\footins}{18pt}	
	
{\linespread{1.25}\selectfont

\vspace*{1.5em}

\begin{center}
%{\relscale{1.2}{\fontfamily{qcr}\fontseries{b}\selectfont 
%{\colorbox{black}{\color{blue}{\llWC{} Database Engine \\and 
%\llMOSAIC{} Native Application Toolkit}}}}}

\colorlet{ctmp}{logoPeach!20!gray}
\colorlet{ctmpp}{ctmp!90!yellow}
\colorlet{ctmppp}{ctmpp!50!black}
\colorlet{ctmpppp}{ctmppp!90!logoRed}
\colorlet{ctmcyan}{ctmpppp!70!cyan}

\colorlet{ctmppppy}{ctmppp!60!orange}

%{\colorbox{darkBlGreen!30!darkRed}{%
\begin{tcolorbox}
[
%%enhanced,
%%frame hidden,
%interior hidden
arc=2pt,outer arc=0pt,
enhanced jigsaw,
width=\textwidth,
colback=ctmppppy!40,
%colback=ctmcyan!50,
colframe=logoRed!30!darkRed,
drop shadow=logoPurple!50!darkRed,
%boxsep=0pt,
%left=0pt,
%right=0pt,
%top=2pt,
]
%\hspace{22pt}
\begin{minipage}{\textwidth}	
\begin{center}	
{\setlength{\fboxsep}{32pt}
	\relscale{1.2}{{\fontfamily{qcr}\fontseries{b}\selectfont%
{The \llMOSAIC{} Publication Repository and 
Structured Research Reporting Framework}
}}}
\end{center}
\end{minipage}
\end{tcolorbox}
\end{center}

\vspace{-3pt}
%\textcolor{darkRed}{\hrule}

\section{Overview}
\p{In recent years --- as publishing has become more \q{digitized} --- there
have emerged many online platforms for scientists to share and publish
their research work, including both raw data and scientific papers.
These portals generally provide an index/search feature where readers
can locate research based on author's name, keywords, or subject
matter, and a document viewer where readers can view each
publication's abstract (and sometimes its full text) and other
publication details (such as co-authors, dates, citations, etc.).
However, these portals offer only limited features for interactively
exploring publications, particularly when it comes to examining data
sets and/or browsing multimedia assets associated with publications,
such as audio, video, interactive diagrams, or 3d graphics.}

\p{The \q{\MOSAIC{}} system is a suite of code libraries allowing 
institutions to host publication repositories which address the 
limitations of existing scientific document portals 
(\MOSAIC{} is an acronym for \q{Multi-Paradigm 
Ontologies for Scientific and Technical Publications}). 
With \MOSAIC{}, each publication can be
linked to a \q{Supplemental Archive} which contains information about
the author's research methods, data sets, and focal topics.  Using 
\MOSAIC{}, developers can implement a hosting platform
for this supplemental material, in addition to publications themselves, 
where the platform provide software
enabling readers to browse and access Supplemental Archives and 
their data sets and/or methodological descriptions.  A \MOSAIC{} 
publication repository, also called a \MOSAIC{} \q{portal,} 
is a structured collection of data and documents which can be 
hosted via web servers (including fully encapsulated and 
containerizable cloud services) implemented with the help of 
\MOSAIC{} libraries.}

\p{\lMOSAIC{} is designed so that the software to access 
publications and Supplemental Archives can be embedded in 
scientific-computing applications, via \q{\MOSAIC{} plugins.}  
This ensures that publications and data sets can be 
examined interactively with the same software as scientists 
use to conduct research.  \lMOSAIC{} also introduces a 
\q{structured reporting system} (\MOSAICSR{}) for describing 
research/experimental/lab methods and protocols.  Supplemental Archives, 
plugins, and \MOSAICSR{} are outlined in the following sections.}

\part{MOSAIC Components}
\section{Supplemental Archives} 
\p{\lMOSAIC{} Supplemental Archives are additional resources 
paired with \MOSAIC{} publications.  In general, Supplemental Archives 
may include raw data, descriptions of research methods, 
or annotated data linked to publication text.  Each Supplemental Archive 
could have any of the following resources:

\begin{enumerate}
\item Interactive versions of the publication, with annotations
indicating important concepts and phrases, perhaps aggregated into a
\q{glossary} defining technical terms;

\item Structured files containing raw data discussed in the publication,
along with interactive software allowing scientists to access and
reuse the data;

\item Detailed reports of the author's research methods and experimental
setup and/or protocols, conforming to the relevant standards with
respect to the publication's subject classification --- for instance,
the \q{Minimum Information for Biological and Biomedical Investigations}
(\MIBBI{}) includes about 40 specific standards for different branches of
biology and medicine;

\item Representations of analytic methods and algorithms used to derive
the research findings, provided directly via computer code or
indirectly via formal descriptions of computational workflows;

\item Self-contained computer software demonstrating code which the
author developed and/or used for analyzing research data;

\item Multi-media assets such as audio or video files, annotated images,
\ThreeD{} graphics, interactive \TwoD/\ThreeD{} plots and diagrams, or other kinds of
non-textual content which needs to be viewed with special multi-media
software.
\end{enumerate}
}

\p{The contents of a Supplemental Archive will be different 
for different publications, depending on whether the archive 
contains specific raw data, or just a summary of the methods 
used to obtain raw data.  In the former case, a typical 
archive will include a \i{data-set application}, or a semi-autonomous 
software component allowing researchers to study and visualize 
the data set, typically provided in turn via data files that are 
also located in the archive.  The data-set application will in 
general provide both a visual interface for raw data and code 
libraries for computationally manipulating this data; typically 
raw data files will encode serialized data structures that 
can be deserialized by code (e.g., \Cpp{} classes) included 
in the data-set application.  \lMOSAIC{} includes generic 
implementations of a \q{data-set explorer} (\MdsX{}) which 
can be adopted to the specific kind of information needing 
to be serialized for a given publication.  The resulting 
data-set application can then be packaged as a plugin to 
existing scientific software, depending on the relevant 
scientific discipline to which the research belongs.}

\section{\protect\llMOSAIC{} Plugins}
\p{By design, \lMOSAIC{} Portals will provide a suite of plugins
for existing scientific software, allowing publications and
Supplemental Archives hosted on the portal to be read and examined
within computer software associated with each publication's subject
matter.  For example, articles about chemistry could be read within
\IQmol{}, a molecular visualization program; articles about cellular
biology and bioimaging could be read within \CaPTk{}, an image-processing
application; articles discussing novel computer code could be read
within \Qt{} Creator, an Integrated Development Environment (\IDE{}) for
programming; etc.  The advantage of accessing
a publication and a Supplemental Archive within actual scientific
software is that it allows research work to be understood, evaluated, 
and reused within the computing environments which scientists actually use to 
conduct professional research.  This is different 
than existing science-publication portals, which generally rely on web browsers to
access supplemental materials like data sets and multimedia files ---
in this setup, as a result, the software ecosystem wherein readers examine
published research is fundamentally separate from the software
platforms where research is actually conducted.  This is one example
of how existing portals are limited in their ability to share research
in rigorously reusable and replicable ways, and how \MOSAIC{} 
offers a potential improvement.}

\p{For authors, embedding presentations of their research within existing
scientific software has the added benefit of making their work 
more practically and immediately useful for the scientific community.  
Such presentations establish the computational framework for practically 
deploying their techniques in real-world scientific contexts, 
accelerating the pace at which research work can be translated to 
concrete scientific (and clinical/lab/experimental) practice.
As one example, research involving novel image analysis techniques 
could be packaged so as to target a \MOSAIC{} plugin for bioimaging
software such as \CaPTk{}, so that readers could actually run the
author's code as a \CaPTk{} module.  This kind of
functional assessment and adoption of novel contributions is harder to carry
out if a research work is only described indirectly within article
text, as opposed to being concretely developed for execution within a
specific scientific application.}

\p{Another benefit of using plugins to 
access Supplemental Archives is that the host application will usually
provide more sophisticated multimedia and data visualization
capabilities, compared to static \PDF{} images or even interactive web
portals.  Publishers have begun to develop online platforms for
browsing research papers in conjunction with multimedia content such
as interactive diagrams and \ThreeD{} graphics --- a physical model of a
protein or a chemical compound, say, can be viewed online via WebGL;
such graphics could even be embedded as an \q{iframe} within an HTML
version of the publications.  Publishers consider this to be
cutting-edge technology.  However, the same molecular \ThreeD{} model, when
viewed in \IQmol{}, can be enhanced with many additional visual features,
representing bonds, orbitals, torsion angles, etc.  The multimedia
experience of exploring chemistry data in custom software like \IQmol{} 
is therefore much greater than the experience of generic web
multimedia, so the scientific software is a better forum for
showcasing novel research.}

\section{\protect\llMOSAIC{} Structured Reporting}

\p{The \MOSAIC{} Structured Reporting framework (\MOSAICSR{}) 
includes tools to help authors develop interactive 
presentations supplementing academic documents, and specifically 
to use Supplemental Archives to document how thier research 
has been conducted.  With \MOSAICSR{}, authors can implement or 
reuse code libraries that report on 
research/experiment methods, workflows, and 
protocols.  The \MOSAICSR{} information may be structured 
as a \q{minimum information checklist} conformant 
to standards such as those collectively 
gathered into the \MIBBI{} recommendations;  
in this case \MOSAICSR{} would be applied by  
implementing object models instantiating \MIBBI{} 
policies.  Alternatively, \MOSAICSR{} reports can be 
derived from actual computer programs simulating 
research workflows, similar to \BioCoder{}\footnote{See 
\bhref{https://jbioleng.biomedcentral.com/articles/10.1186/1754-1611-4-13}.} (which 
is included, in an updated \Cpp{} version, as one 
\MOSAIC{} library).  Finally, \MOSAICSR{} presentations 
may be based on annotations applied to research/analytic 
code.  For example, in the context of image analysis, 
the \Pandore{} project provides an \q{Image Processing 
Objectives} Ontology as well as a suite of image-processing 
operations that can be called from computer code.  
Image-analysis pipelines can therefore be explained 
by annotating the pertinent function-calls 
(which \Pandore{} calls \q{operators}) with terms from 
the \Pandore{} controlled vocabulary.}

\p{\lMOSAICSR{} can express both computational 
workflows that are fully encapsulated by published 
code and real-world protocols concerning 
laboratory equipment and physical materials 
or samples under investigation.  In the latter 
guise, \MOSAICSR{} code can employ or instantiate 
standardized terminologies and data structures 
for describing experiments --- such as 
\MIBBI{} policies or \BioCoder{} functions.  
In this case, the role of \MOSAICSR{} code is 
to serve as a serialization/deserialization 
endpoint for sharing research metadata.  
Conversely, when workflows are fully implemented 
within software developed as part of a body 
of research, \MOSAICSR{} can provide a functional 
interface allowing this code to be embedded 
in scientific software.  For these latter cases, 
\MOSAICSR{} provides a framework for modeling 
how a software component specific to a given 
research project exposes its functionality to 
host and/or networked peer applications.  There 
are also cases where both scenarios are 
relevant --- the \MOSAICSR{} code would simultaneously 
document real-world experimental protocols and 
construct a digital interface as part of 
a workflow which is part digital (\q{\textit{in silico}}) and part 
\q{real-world} (\q{in the lab}).}

\section{\protect\llMOSAIC{} Annotations}
\p{Included in any \MOSAIC{} plugin would be specially-designed \PDF{} viewers for
interactively reading student's papers.  In particular, these \PDF{}
applications would recognize cross-references linking between student
publications and their associated Supplemental Archives.  This would
allow readers to identify concepts which are discussed and/or
represented both in the research paper and in the archive, annotating
their papers accordingly.  For example (a similar example is used in
the documentation for \BioCoder{}, the concept \q{RNA Extraction} may be discussed in a
publication text, and also formally declared as one step in the lab
processing as  represented via \BioCoder{}, and summarized in a
\BioCoder{}-generated chart (and included in the Supplemental Archive).
The \PDF{} viewer would then ensure that the phrase \q{RNA Extraction} in
the text is interactively linked to the concordant step in the
experimental process, so that readers could view the \BioCoder{} summary
as a context-menu action associated with the phrase where it appears
in the \PDF{} file.  For a different example, the phrase \q{Oxygenated
Airflow} refers to airflow in assisted-breathing devices, such as
ventilators; to ensure that the device is working properly, the
equipment must be monitored to ensure that a steady stream of oxygen
reaches the patient.  Research into the design and manufacturing of
ventilators and similar devices may then include \q{Oxygenated Airflow}
both as a phrase within the article text and as a parameter in data
sets evaluating the device's performance.  In this situation, the
publication location of the \q{Oxygenated Airflow} phrase should once
again be annotated with links to the relevant part of the data set
(e.g., a table column) where measurements of airflow levels are
recorded.}

\part{Concrete Applications}
\section{\protect\llMOSAIC{} in the context of Image Analysis}
\p{This section will focus on one specific application 
of \MOSAICSR{} in the context of image analysis 
and bioimaging --- specifically, a \q{Data Structure 
Protocol for Image-Analysis Networking} (\DSPIN{}), 
which both extends and adds a narrower focus 
to \MOSAICSR{}.  Software using the \DSPIN{} 
protocol provides a description of image 
processing capabilities which have been 
utilized and/or are functionally exposed 
by code and data associated with a research 
project.  This includes \q{structured reporting} 
of research objectives as well as a concrete 
interface for invoking analyses associated 
with the research (either new algorithms or 
techniques used to obtain reported findings).
\lDSPIN{}, in turn, is based on \CaPTk{} 
(the Cancer Imaging and Phenomics Toolkit) 
and \Pandore{} (an image-processing environment 
which includes both data models and interactive 
software).  The \Pandore{} \q{Image Processing Objectives} 
Ontology, mentioned above,  
that provides a structural basis for \DSPIN{}.  
For information about how different objectives 
are merged into workflows, \DSPIN{} 
adopts protocols from \CaPTk{}, particularly 
with respect to implementating 
image-analysis capabilities as 
extensions to a core application, and 
with respect to \CaPTk{}'s implementation of the 
Common Workflow Language (\CWL{}).  In effect, 
\DSPIN{} formalizes the data models and 
prototypes adopted by \Pandore{} and 
\CaPTk{} so as to concretize \MOSAICSR{} 
for the specific domain of image processing 
and Computer Vision.  The following sections 
will therefore outline \DSPIN{} features 
in the context of \MOSAICSR{} design principles 
and objectives.}

\section{Meta-Procedural Modeling in 
\protect\lsDSPIN{} and \protect\lsMOSAICSR{}}

\p{Most approaches to modeling research workflows 
involve some concept of \q{meta-objects},\footnote{See 
the \sVISSION{} system: \bhref{https://pdfs.semanticscholar.org/1ad7/c459dc4f89f87719af1d7a6f30e6f58dff17.pdf}.} 
\q{tools} (in the terminology of \CWL{}), 
or \q{transitions} (in the language of Petri Net 
theory).  In \MOSAICSR{}, the analogous concept is 
that of \textit{meta-procedures}, which are analogous 
to ordinary computational procedures but 
add extra sources of information concerning 
input and output parameters.  In general, 
rather than simply passing an imput value 
into an executable routine, metaprocedures 
define steps which can be taken to 
acquire the proper values when needed.  
Aside from ordinary runtime values, the 
most important input sources are methods 
defined on \GUI{} components; command-line 
parameters; file contents; and 
not-yet-evaluated expressions (perhaps 
encapsulated in scripts or function pointers).  
A meta-procedure formulation abstracts the 
acquisition of input arguments (or the consumption of 
values within \q{channels} via which a procedure sends and 
receives data) 
from the concrete procedure or procedures 
which are eventually executed.  Therefore, 
a \MOSAIC{} meta-procedure definition 
has two separate parts: a preamble where 
input sources are described; and an 
executive sequence where concrete procedures 
are indicated.  A \textit{meta-evaluator} 
then operates in accord with these 
definitions, concretizing the input values 
and launching the actual procedure(s).  
For \DSPIN{}, meta-procedures can be defined 
using a framework based on \BioCoder{}, 
but adopted to the imaging and Computer Vision 
context.}

\p{Image analysis methods are often 
described in academic literature in 
terms of mathematical formulae and/or 
characterizations of computing environments 
(such as Graphical Processing Units); 
it requires additional construction to 
translate these overviews into actual 
computer code.  Once 
Computer Vision innovations are in fact 
concretely implemented, there is then 
an additional stage of development 
requisite for users to actually enact 
the computations described in the 
research.  Although it is theoretically 
possible to demonstrate novel methods 
within fully self-contained autonomous 
applications, it is more convenient for 
users if research code is integrated 
with existing imaging software.  
Along these lines, the \DSPIN{} interface can help 
connect new code to existing applications, 
allowing users to access the new code's functionality 
through \GUI{} actions, command-line invocations, 
or inter-application messaging.}

\p{In addition to practically enabling 
application embedding, \DSPIN{} models 
represent research methods and theories, 
contributing to transparency and 
reusability according to the 
\MIBBI{} and \FAIR{} (Findable, Accessible, 
Interoperable, Reusable) standards.\footnote{See 
\bhref{https://www.researchgate.net/publication/331775411_FAIRness_in_Biomedical_Data_Discovery}.}
This can be achieved, in part, by implementing 
data structures conforming to \Pandore{} 
Image Processing Objectives.  However, \DSPIN{} 
embeds this logic in an Object-Oriented 
context which allows imaging-specific workflow 
notations to be paired with specifications 
outside of image-processing in the narrowest 
sense.  This allows \DSPIN{} to be available 
for hybrid computational objectives representations 
which are only partially covered by the imaging 
domain --- analogous to the \MIAPEGI{} 
(Gel Electrophoresis Informatics) component 
of \MIAPE{} (Minimum Information About a Proteomics 
Experiment).  The following section will 
discuss several domains where \DSPIN{} has 
been explicitly integrated with code 
libraries codifying \MIBBI{}-style research protocols.}

%\vspace{-1em}
\section{\protect\lsDSPIN{} in Contexts Supplemental to Image Processing.}
\vspace{2em}
\subsection{Image Flow Cytometry}
\p{One important use-case for biomedical image processing 
is to analyze cellular microscopy in conjunction with 
cytometric experiments which investigate cells and 
cellular-scale entities (such as proteins) indirectly.  
Conceptually, image analysis and flow cytometry 
(\FCM{}) analysis 
are mathematically similar, and some commercial 
cytometry software has been extended with 
image-processing capabilities.  The overlap 
between cytometric and image analysis has also 
inspired attempts to merge cytometry standards, 
such as \MIFlowCyt{} (the Minimum Information about a Flow Cytometry Experiment policy within \MIBBI{}), with bioimaging 
standards such as \DICOM{} (Digital Imaging and 
Communications in Medicine).  One such proposal is 
due to Robert Leif, who argues that \q{The large overlap between imaging and flow cytometry provides strong evidence that both modalities should be covered by the same standard} 
and has formalized an \XML{} language 
(\CytometryML{}) to serve as that 
overarching bridge.\footnote{See 
\bhref{https://spie.org/Publications/Proceedings/Paper/10.1117/12.2295220?SSO=1}.}  The \DSPIN{} project builds 
off this work by introducing its own \FCM{}/\DICOM{} 
hybrid, although in an object-oriented rather 
than \XML{}-based context (along with an expanded 
\XML{}-oriented schema discussed in 
the next section).  As a reference implementation 
for this \DSPIN{} extension, the project 
also provides a pure-\Cpp{} cytometry 
library based on the \openCyto{} and 
\FACSanadu{} libraries, but eliminating 
external dependencies such as \R{} and 
\Java{}.  The \FCM{}/\DICOM{} bridge is 
implemented in this context via a 
\DSPIN{} supplement to \DICOM{} based 
on \q{Semantic \DICOM{},} which is an 
effort to standardize query processing 
within \PACS{} (Picture Archiving and 
Communications Service) workstations and 
to more effectively integrate \DICOM{} with 
clinical data.}

\subsection{A Semantic \protect\lsDICOM{} Object Model}

\p{As a formal representation of imaging workflows, 
\DSPIN{} would reasonably be paired in 
many contexts with \DICOM{}, insofar as \DICOM{} 
represents the canonical standard for exchanging 
medical image data.  For its applications within 
the medical-imaging context, therefore, 
\DSPIN{} provides object-oriented accessors 
to \DICOM{} data such that image-objectives and 
\DICOM{} object models can interoperate.  
This object-oriented foundation also provides a 
basis on which to further integrate clinical 
data in the form of \q{Semantic} \PACS{} 
models.\footnote{See \bhref{https://www.ncbi.nlm.nih.gov/pmc/articles/PMC5119276/}.}  The original Semantic \PACS{} implementation 
draws clinical data from \DICOM{} headers, and 
represents this extracted information via \RDF{}.  
In keeping with \MOSAIC{}'s more object-oriented 
focus, the \MOSAIC{}-Semantic \DICOM{} (\MOSAICSD{}) 
integration is engineered instead as an 
extension to the \DICOM{} Toolkit (\DCMTK{}) 
library, though it imports many constructs 
within Semantic \DICOM{} in the guise 
of a \q{Hypergraph Ontology}, one example of a 
system \MOSAIC{} uses to merge Semantic Web 
schema with Object-Oriented code.  In 
short, \MOSAICSD{} 
extends \DSPIN{} by defining  
a central \DICOM{} object model affixed 
to both clinical and 
image-processing object models.}

\p{The \MOSAICSD{} object system applies 
not only to \DICOM{} integrated with 
statements of image-processing objectives, 
but also to other biomedical contexts where 
image analysis should be integrated 
with other analytic modalities and 
also with clinical or epidemiological 
information.  For example, Flow 
Cytometry overlaps with clinical data tracking because 
one of \FCM{}'s essential investigative roles 
is to examine patients' immunological 
response to diseases and/or interventions.  
In the context of Covid-19, say --- with respect to 
achieving a deeper understanding of how 
and why SARS-CoV-2 symptoms present differently 
in different patients --- \q{The starting point will likely be a deep characterization of the immune system in patients with different stage of the disease}.\footnote{See \bhref{https://onlinelibrary.wiley.com/doi/full/10.1002/cyto.a.24002}.}  
That is, \FCM{} observations need to be 
matched with clinical data in order to 
classify (and consider statistical correlations 
between) immunological findings and clinical facts 
(risk factors, sociodemographics, disease 
progression, and so forth).  This analysis 
intrinsically assumes that \FCM{} data can 
be transparently linked with all relevant 
clinical data, but such integration is 
difficult even in \DICOM{}, where \DICOM{} 
headers are specifically designed for preserving 
patient data across picture-sharing networks  
(there is no analogous \q{header} 
component in the 
Flow Cytometry Standard, \FCS{}).  
In brief, \DSPIN{} extensions to 
non-imaging domains can promote 
data integration insofar as \DSPIN{} 
Clinical Object Models, based on 
Semantic \DICOM{}, provide 
an affixation point for clinical data 
in analytic contexts which are 
operationally related to image analysis, 
not just to image analysis per se.}

\subsection{Geo-Imaging and Geographic Information Systems}
\p{Another area where \DSPIN{} provides structured 
object models is that of Geographic Information Systems 
(\GIS{}) annotations.  
There is a direct link between image processing 
and \GIS{} insofar as identifying geotaggable 
features is one dimension or application of 
Computer Vision.  Effectively manipulating 
geoimaging data requires mathematical 
translations between several different 
coordinate systems, in both two and 
three dimensions.  These coordinate 
transforms --- as well as semantic interpretations 
of geoimage segments (buildings, land features, 
roads, etc.) --- can serve as the basis for 
an object model attaching image-processing 
objectives to \GIS{} workflows.}

\p{In conventional \GIS{} annotation, 
data structures are linked to both 
geospatial coordinates and to 
visual cues or icons  
allowing locations of interest to be 
indicated on maps.  The actual geotagged 
data structures can be derived from 
any domain; as such any object model 
can be integrated with \GIS{} annotations 
so long as one can assign spatial interpretations 
to the phenomena computationally encapsulated 
by the domain in question.  In a medical 
context, for instance, geotagged data 
may represent the scope of a vaccination campaign 
or the extent of an epidemic, along with 
relevant geographic or civic features 
(villages, medical clinics, national borders, 
and so on).  The actual map as a visible 
digital artifact therefore serves as a 
virtual glue where clinical, geographical, 
governmental, and \GUI{} data are all 
sutured together.  Insofar as geoimaging 
involves analysis of photographed land 
features and/or urban environments, 
image processing information represents a 
further object model that can be 
added to the mix.  Even when dealing solely 
with virtual maps, however --- rather 
than with satellite images or other geospatial 
photography --- the analysis and application-level 
rendering of map features is sufficiently 
similar to image processing that a rigorous 
model of \GIS{} integration belongs 
properly within \DSPIN{}, specifically 
within a \MOSAICGIS{} extension.}

\part{MOSAIC Data Models}

\section{Conceptual Space Models and Hypergraph Database Implementation}

\p{}


\section{Markup Serialization and \protect\q{Grounding}}

\p{The data-integration mechanisms discussed 
in this paper thus far have focused on 
object models and object-oriented programming 
techniques.  While composing 
special-purpose object-based libraries 
specifically tailored to individual  
data-integration problems is a powerful 
tool for solving such problems, data integration initiatives 
in practice are 
often organized around standards 
for sharing or serializing conformant 
data structures.  As a result, an important 
aspect of data integration is implementing 
proper data-serialization technology 
which is sufficiently rigorous to 
serve as a proxy for formal interface 
definitions.  
\lXML{} formats are a good case-study 
for such formalizations because 
most serialization in 
the biomedical and bioimaging context, 
operates through 
\XML{} languages.}

\p{Standardizing a data 
format by stipulating how \XML{} files 
may encode the data is simpler than 
defining an analogous specification 
in terms of executable computer code: 
one way to document the shape of any 
relevant data is just to explain how an \XML{} 
document will be structured insofar as 
it encodes data accordingly.  Potentially, 
such specification can be a single 
\XML{} \DTD{} file, or an \XML{} sample, 
providing a convenient reference point 
for developers to grasp the underlying 
data model.  However, the structure of an \XML{} document 
does not, in itself, present a clear picture of 
how the information which the document 
represents is semantically organized.  Even 
though \XML{} is processed by computer programs, 
it is not even evident from an \XML{} document or 
schema which \XML{} elements (if any) correspond 
to data types recognized by applications which 
read and/or write the corresponding \XML{} code.  
For example, the \XML{} portion of \OMETIFF{} 
(the principal Open Microscopy Environment imaging 
format) includes an explicit \textbf{Image} element 
(which gathers up all significant image metadata); 
an application reading \OMETIFF{} files 
might therefore introduce a single datatype 
--- analogous to (and maybe wrapping an) \textbf{itk::Image} from 
the Insight Toolkit imaging libraries --- 
bundling the data in that part of the \OME{} \XML{}.  
In this case there will be a one-to-one correspondence 
between \XML{} structure and application-level 
data types, at least for that one \textbf{Image} node.  
On the other hand, software reading \OMETIFF{} 
information may not manipulate images directly, 
but rather pull out other kinds of metadata, 
such as an experimenter's name or description 
of the microscopic setup.  In this case, 
the application may not have an explicit 
\q{object} representing the image itself, 
but it may still read information about the 
image from child nodes of the \XML{} \textbf{Image} 
element.}

\p{In short, applications can read or use data 
from an \XML{} document in different ways, 
so the document's structure does not itself 
provide a clear picture of how code which 
reads the \XML{} is organized.  This uncertainty 
is significant insofar as one wishes to use 
\XML{} specifications as an indirect strategy 
for documenting parameters and features 
of the data structures which are serialized 
via the relevant \XML{} language.  In effect, 
\XML{} serialization operates on two levels: 
on the one hand, the specific \XML{} document 
provides an encoding of data conforming to a 
given structure; but, at a more abstract 
level, one can model the relationship 
between the surface-level \XML{} node 
structure and the application-level data structures 
thereby serialized.  This second level of detail 
is usually implicit and unstated, but in \MOSAIC{}, such 
\q{meta-serialization} --- a term used 
to suggest the idea of providing meta-data 
\textit{about} a serialization --- is directly 
formulated through a notion of 
\q{grounding}, which involves adding supplemental 
markup clarifying how documents instantiating a 
serialization format (such as \XML{}) relate 
to the coding protocols and data types of 
software which reads or creates these documents.}

\p{For maximum expressivity, \MOSAIC{} introduces a 
meta-serialization system that can be applied 
to languages more flexible than \XML{} --- notably 
\TAGML{} --- as well as to \XML{} 
proper.\footnote{See \bhref{http://www.balisage.net/Proceedings/vol21/print/HaentjensDekker01/BalisageVol21-HaentjensDekker01.html}.}
In effect, \MOSAIC{} provides parsers for an 
extended version of \TAGML{} which includes 
an additional \q{grounding} layer.  Grounding, 
in this context, means describing how 
elements in the markip --- nodes, attributes, 
and character sequences --- are \q{grounded} 
in application-level types, data fields, 
and other programming constructs.  \lMOSAIC{} 
provides parsers for this \q{Grounded \TAGML{}} 
(\GTagML{}) language as well as converters 
to output serialized data as conventional 
\XML{}, so that \GTagML{} specifications can 
be used in contexts (\MIFlowCyt{}, for instance) 
where ordinary \XML{} is expected and serves 
as a basis of standardization.  \lMOSAIC{} also 
provides code libraries 
for extracting data from \GTagML{} documents.}

\p{\lGTagML{} documents, 
with certain restrictions enforced, are structurally 
isomorphic to \XML{} and can be rendered as 
pure \XML{}; as such, \GTagML{} schemata can 
be used to define norms for \XML{} languages, 
although the logical role and operations of such requirements 
is not identical to \XML{} schema definitions.  
By design, \GTagML{} schemata operate on 
two levels: on the one hand, they constrain 
the organization of \GTagML{} files themselves; 
but they also stipulate how \GTagML{} document 
structure relates to the type systems of 
code libraries serializing and deserializing 
\GTagML{} files.  To model this second, type-oriented 
metadata, \GTagML{} introduces a hypergraph-based 
type model organized around \q{infosets,} which 
are structured overviews of application-level 
data types.  To fully utilize \GTagML{} features, 
programmers may then compose \q{infoset classes} 
as wrappers around ordinary data types 
(e.g., \Cpp{} classes) whose instances are 
serialized via nodes or character strings 
within \GTagML{} documents.  Infoset classes 
provide hypergraph \q{views} onto type-instances, 
and act as a bridge between data types and 
their associated \GTagML{} schema.  In particular, 
infoset classes (rather than the types 
they encapsulate) are the basis for schematizing 
the relation between \GTagML{} document elements 
and type instances (and the data they contain).}

\p{The hypergraph-based modeling incorporated 
into \GTagML{} reuses much of the code associated 
with \ConceptsDB{}, which is a new hypergraph database 
being developed alongside \MOSAIC{}.  More information 
about \ConceptsDB{} can be found on 
the guthub project page for Linguistic 
Technology Systems (\bhref{https://github.com/Mosaic-DigammaDB/LingTechSys}).}



\end{document}


